% ============================================================
%  Cl(6) Lattice QCD — Paper 2
%  Todd M. Firestone
%  DRAFT — not for circulation
%  Target journal: Advances in Applied Clifford Algebras
%  Typeset for arXiv submission
%  Document class: article + geometry (generic, arXiv-compatible)
%
%  FLAGS USED IN THIS FILE:
%    % NOTE:   — placeholder or missing content
%    % REVIEW: — internal inconsistency; author decision required
%    % ASK:    — ambiguity that cannot be resolved without author input
% ============================================================

\documentclass[12pt]{article}

% --- Encoding & fonts ---
\usepackage[utf8]{inputenc}
\usepackage[T1]{fontenc}
\usepackage{lmodern}
\usepackage{microtype}

% --- Page geometry (arXiv standard margins) ---
\usepackage[top=1in, bottom=1in, left=1.25in, right=1.25in]{geometry}

% --- Mathematics ---
\usepackage{amsmath}
\usepackage{amssymb}
\usepackage{amsthm}
\usepackage{mathtools}

% --- Tables ---
\usepackage{booktabs}
\usepackage{array}
\usepackage{multirow}
\usepackage{longtable}

% --- Graphics & color ---
\usepackage{graphicx}
\usepackage{xcolor}

% --- Citations: numbered, sorted, compressed [1,2,3] style ---
\usepackage[numbers,sort&compress]{natbib}
\bibliographystyle{unsrtnat}

% --- Hyperlinks (arXiv-safe) ---
\usepackage[
colorlinks=true,
linkcolor=blue!60!black,
citecolor=blue!60!black,
urlcolor=blue!60!black,
pdftitle={Grade Decomposition of the Lattice QCD Pion Correlator
	in the Cl(6) Framework},
pdfauthor={Todd M. Firestone}
]{hyperref}

% --- Equation numbering: sequential (1),(2),(3)... ---
% (arXiv single-paper standard; no section prefix)

% ============================================================
%  THEOREM ENVIRONMENTS
% ============================================================
\newtheorem{theorem}{Theorem}[section]
\newtheorem{lemma}[theorem]{Lemma}
\newtheorem{proposition}[theorem]{Proposition}
\newtheorem{corollary}[theorem]{Corollary}
\theoremstyle{definition}
\newtheorem{definition}[theorem]{Definition}
\theoremstyle{remark}
\newtheorem{remark}[theorem]{Remark}

% ============================================================
%  CUSTOM MACROS
% ============================================================

% --- Inherited from Paper 1 (unchanged) ---
\newcommand{\cl}{\mathrm{Cl}}           % Clifford algebra: \cl(6)
\newcommand{\spin}{\mathrm{Spin}}       % Spin group: \spin(6)
\newcommand{\su}{\mathrm{SU}}           % SU group: \su(3)
\newcommand{\tr}{\mathrm{Tr}}           % Trace
\newcommand{\dW}{D_{\mathrm{W}}}        % Wilson--Dirac operator
\newcommand{\SW}{S_{\mathrm{W}}}        % Wilson action
\newcommand{\kc}{\kappa_{\mathrm{c}}}   % critical hopping parameter
\newcommand{\half}{\tfrac{1}{2}}

% --- New in Paper 2 ---
\newcommand{\kcren}{\kappa_{c}^{\mathrm{ren}}}      % renormalized kappa_c
\newcommand{\GaSeven}{\Gamma_7}                      % Spin(6) chirality
\newcommand{\Isix}{I_6}                              % Cl(6) pseudoscalar
\newcommand{\Qtil}{\widetilde{Q}}                    % 24x24 matrix in Theorem~\ref{thm:kinematic}
\newcommand{\mPCAC}{m_{\mathrm{PCAC}}}               % PCAC mass
\newcommand{\ftriplet}{f_{\mathrm{triplet}}}          % combined triplet fraction
\newcommand{\fsinglet}{f_{\mathrm{singlet}}}          % combined singlet fraction
\newcommand{\dftriplet}{\delta f_{\mathrm{triplet}}}  % triplet phase difference
\newcommand{\Ncfg}{N_{\mathrm{cfg}}}                  % number of configurations
\newcommand{\Ntherm}{N_{\mathrm{therm}}}              % thermalization sweeps

% ============================================================
%  DRAFT HEADER (matches Paper 1)
% ============================================================
\usepackage{fancyhdr}
\pagestyle{fancy}
\fancyhf{}
\rhead{\textit{Cl(6) Lattice QCD $\cdot$ DRAFT --- not for circulation}}
\rfoot{\thepage}
\renewcommand{\headrulewidth}{0.4pt}

% ============================================================
\begin{document}
	% ============================================================
	
	% --- Title block ---
	\title{%
		\textbf{Grade Decomposition of the Lattice QCD Pion Correlator}\\[4pt]
		\textbf{in the \texorpdfstring{$\cl(6)$}{Cl(6)} Framework}%
	}
	
	\author{Todd M.\ Firestone}
	
	% NOTE: Replace [Date] with actual submission date before arXiv/journal submission.
	\date{%
		\small\textit{ }%
	}
	
	\maketitle
	\thispagestyle{fancy}
	
	% ============================================================
	\begin{abstract}
		% ============================================================
		The Clifford algebra $\cl(6)$ embeds the $\su(3)$ color gauge group as a
		subgroup of $\spin(6)\cong\su(4)$, making the $\su(3)$ representation
		structure of the Wilson--Dirac spinor space algebraically explicit through
		the decomposition
		$\mathbf{8} = \mathbf{3}\oplus\bar{\mathbf{3}}\oplus\mathbf{1}\oplus\mathbf{1}$.
		We construct the corresponding four-channel decomposition of the lattice
		pion correlator and report its first numerical implementation on quenched
		gauge ensembles.
		
		The four representation projectors are constructed as Lagrange spectral
		projectors of the Cartan hypercharge $H = B_{12}+B_{34}+B_{56}$ of
		$\mathrm{spin}(6)$, whose eigenspectrum
		$\{\pm\tfrac{1}{2},\,\pm\tfrac{3}{2}\}$ matches the $\su(3)$
		representation dimensions.  They satisfy completeness, idempotency,
		mutual orthogonality, $\su(3)$ covariance, charge-conjugation pairing, and
		commutativity with the 4D chiral operator $\gamma_5$ to numerical residual
		$\leq 2.4\times 10^{-15}$.  The labeling $(\GaSeven, H)$, where
		$\GaSeven = i\Isix$ is the $\spin(6)$ chirality, is the standard
		$\su(4)\supset\su(3)\times\mathrm{U}(1)$ Levi decomposition of the Dirac
		spinor of $\spin(6)$, known from representation theory and from the
		Pati--Salam $\su(4)_{\mathrm{color}}$ model.
		
		A kinematic limit theorem follows from the algebraic structure: in the
		free Wilson--Dirac theory, the grade fractions
		$f_r(t)\equiv\langle C_\pi^{(r)}(t)\rangle/\langle C_\pi(t)\rangle$
		equal $\dim(r)/8$ exactly, a consequence of
		the residual $I_4{\otimes}U(2)$ symmetry of the free operator combined
		with $\su(3)$ color and spatial-cubic invariance.
		
		On thermalized quenched ensembles at $N=6$, $\beta\in\{4.5,6.0\}$, and
		bare quark mass $m_0=0.05$, the algebraic completeness identity
		$\sum_r C_\pi^{(r)}(t)=C_\pi(t)$ is verified per configuration to
		maximum residual $5.4\times 10^{-16}$, and the $\spin(6)$-chirality fractions
		are fixed at $f_{(\GaSeven=\pm1)}=\tfrac12$ exactly on every configuration by
		an identity independent of the gauge dynamics.  The measured grade fractions are
		consistent with the kinematic free-field values at the $10^{-4}$ level,
		and the phase differences between confined and deconfined regimes are
		consistent with zero at heavy quark mass.  The decomposition admits no
		natural analog within the standard four-component Wilson--Dirac algebra:
		the bivector $B_{56}$ requires the internal generators $e_5,e_6$ that
		are absent from $\cl(4)$, and the resulting grade structure cuts across
		the $\gamma_5$ chirality decomposition.  The chiral-regime extension at
		$\kappa\to\kcren$, required to test whether non-perturbative gauge
		dynamics distinguish the four representation channels, is identified as
		the target for follow-up investigation.
	\end{abstract}
	
	\bigskip
	\noindent\textit{Keywords:} Clifford algebra, lattice QCD, Wilson
	fermions, pion correlator, grade decomposition, Cartan hypercharge,
	$\su(3)$ representation theory, kinematic limit theorem, quenched
	approximation.
	
	% ============================================================
	\section{Introduction}
	\label{sec:intro}
	% ============================================================
	
	\subsection{Motivation}
	\label{sec:motivation}
	
	The Clifford algebra $\cl(6)$ provides a natural home for the internal
	algebraic structure of Wilson lattice QCD.  Its eight-dimensional minimal
	left ideal decomposes under $\su(3)$ as
	$\mathbf{3}\oplus\bar{\mathbf{3}}\oplus\mathbf{1}\oplus\mathbf{1}$---accommodating
	quark color triplets, antiquark anti-triplets, and two gauge singlets---while
	the 15 bivectors span the Lie algebra of $\spin(6)\cong\su(4)$, which
	contains $\su(3)$ as the stabilizer subgroup of a preferred spinor direction.
	This algebraic architecture was established by
	Furey~\cite{Furey2012,Furey2014,Furey2015,Furey2018a,Furey2018b} and is
	reviewed in Paper~1~\cite{FirestoneP1}; what was missing prior to Paper~1
	was any dynamical calculation within the framework.
	
	Paper~1~\cite{FirestoneP1} supplied that missing step: a complete,
	benchmark-verified Wilson lattice QCD simulation was constructed within
	$\cl(6)$, demonstrating numerical equivalence with standard Wilson lattice
	QCD across gauge and fermionic observables.  Section~5.6.4 of that paper
	gestured toward what is now the subject of the present work:
	\emph{grade-decomposed observables}, defined in terms of the $\cl(6)$ grade
	structure and carrying $\su(3)$ representation-theoretic content that the
	standard four-component formulation cannot express.
	
	This paper delivers what Section~5.6.4 of Paper~1 gestured at.  We construct
	the grade decomposition of the pion correlator explicitly, prove a kinematic
	limit theorem that governs its free-field behavior, and implement it on Monte
	Carlo gauge ensembles.  The motivating physical question is whether
	non-perturbative gauge dynamics---confinement, chiral symmetry breaking, the
	quantum vacuum---distribute spinor weight unequally among the four
	representation channels.  The present work establishes the algebraic machinery
	and confirms the null result at heavy quark mass.  The chiral-regime test of
	the central question is identified as future work.
	
	\subsection{Prior Work and the Gap}
	\label{sec:priorwork}
	
	Three distinct bodies of literature are relevant to the present paper, and
	none overlaps with its essential contribution.
	
	The first is the algebraic classification program in $\cl(6)$ and related
	division algebras~\cite{Furey2012,Furey2014,Furey2015,Furey2018a,Furey2018b,
		Dixon1994,Stoica2018}.  Furey's
	work~\cite{Furey2012,Furey2014,Furey2015,Furey2018a,Furey2018b} identifies
	the quantum numbers of a single Standard Model generation from the action of
	$\cl(6)$ on its minimal left ideal, and connects the
	$\su(3)\times\su(2)\times\mathrm{U}(1)$ symmetry to the algebra's automorphism
	structure.  Stoica~\cite{Stoica2018} extends related constructions.
	Dixon~\cite{Dixon1994} works in the division algebra context.  None of these
	works computes dynamics, and none formulates a lattice theory.
	
	The second is the Pati--Salam $\su(4)_{\mathrm{color}}$
	model~\cite{PatiSalam1974} and the associated branching rules of the
	$\su(4)\supset\su(3)\times\mathrm{U}(1)$ Levi decomposition, tabulated by
	Slansky~\cite{Slansky1981}.  These establish the representation-theoretic
	content of the $(\GaSeven, H)$ double grading at the group theory level; the
	present paper is the first work to implement this grading dynamically.
	
	The third is the broader landscape of alternative representations for
	$\su(3)$ lattice gauge theory---loop-string-hadron formulations, tensor network
	encodings, qubit-compatible embeddings; see Paper~1~\cite{FirestoneP1}~\S1.2
	for a full survey.  None of these uses the Clifford algebra as the
	reformulation language, and none introduces grade-projected correlators.
	
	The gap filled by this paper is the first dynamical implementation of the
	$(\GaSeven, H)$ grading: construction of the projectors $P_r$, proof of the
	kinematic limit theorem, and measurement of the grade fractions on thermalized
	gauge ensembles.
	
	\subsection{Scope and Organization}
	\label{sec:scope}
	
	This paper makes three contributions.  First, the algebraic construction of
	the four-channel grade decomposition of the pion correlator, including the
	Cartan hypercharge $H$ as the grading operator and the complete set of
	projector identities (Secs.~\ref{sec:framework}
	and~\ref{sec:projectors}).  Second, a kinematic limit theorem proving that
	free-field grade fractions equal $\dim(r)/8$ exactly, with a full algebraic
	proof from the $I_4{\otimes}U(2)$ symmetry of the free Wilson--Dirac operator
	(Sec.~\ref{sec:kinematic-theorem}).  Third, the first numerical implementation
	on quenched $\su(3)$ gauge ensembles, with $N=6$, $\beta\in\{4.5,6.0\}$, and
	$m_0=0.05$ (Sec.~\ref{sec:experiment-A}).
	
	This paper does not claim to have measured a dynamical grade-fractionation
	signal.  The heavy-mass data is consistent with the kinematic theorem at the
	$10^{-4}$ level; the chiral-regime extension needed to test the central
	hypothesis is identified as Experiment~B
	(Sec.~\ref{sec:outlook}).
	
	The present work is deliberately restricted to the grade-decomposition
	observable.  A broader programme---renormalized Gell-Mann--Oakes--Renner
	analysis and Banks--Casher spectral measurements---was considered but is
	not pursued here: the renormalized-GOR signal is dominated by additive
	Wilson mass renormalization and does not cleanly isolate the physics of
	interest, and the spectral-density measurement belongs to the
	chiral-regime extension (Experiment~B, Sec.~\ref{sec:experiment-B-design}).
	The $\kcren$ determination of Sec.~\ref{sec:kappa-c} is retained solely as
	a prerequisite for that extension, not as input to any renormalized-GOR fit.
	
	Three limitations govern the current results:
	single lattice volume ($N=6$), single lattice spacing ($\beta\in\{4.5,6.0\}$
	fixed), and quenched approximation throughout.
	
	\medskip
	\noindent\textit{Conventions.}\enspace
	Euclidean signature throughout.  Spacetime indices $\mu,\nu\in\{1,2,3,4\}$;
	$\cl(6)$ internal indices $i,j\in\{1,\ldots,6\}$.  Wilson parameter $r=1$
	fixed; $\kappa=1/[2(m_0+4r)]$.  All conventions match
	Paper~1~\cite{FirestoneP1} exactly; that paper is the primary foundational
	reference for algebraic details not reprised here.
	
	% ============================================================
	\section{Algebraic Framework}
	\label{sec:framework}
	% ============================================================
	
	\subsection{\texorpdfstring{$\cl(6)$}{Cl(6)} Basis and Grade Structure}
	\label{sec:basis}
	
	The Clifford algebra $\cl(6)$ is generated by six orthonormal basis vectors
	$\{e_1,\ldots,e_6\}$ satisfying
	\begin{equation}
		e_i e_j + e_j e_i = 2\delta_{ij}, \qquad i,j\in\{1,\ldots,6\}.
		\label{eq:anticomm}
	\end{equation}
	An $8\times8$ matrix representation is given by Pauli tensor products:
	\begin{equation}
		\begin{aligned}
			e_1 &= \sigma_x\otimes I_2\otimes I_2, &\quad
			e_2 &= \sigma_y\otimes I_2\otimes I_2, \\
			e_3 &= \sigma_z\otimes\sigma_x\otimes I_2, &\quad
			e_4 &= \sigma_z\otimes\sigma_y\otimes I_2, \\
			e_5 &= \sigma_z\otimes\sigma_z\otimes\sigma_x, &\quad
			e_6 &= \sigma_z\otimes\sigma_z\otimes\sigma_y.
		\end{aligned}
		\label{eq:basis}
	\end{equation}
	All six are Hermitian and square to $I_8$; the anti-commutation
	relations~\eqref{eq:anticomm} hold to maximum error $0.00\times10^0$ in
	IEEE~754 double precision.  We refer to Paper~1~\cite{FirestoneP1}
	Appendix~A for the complete 64-element basis enumeration and multiplication
	table.  The grade decomposition
	\begin{equation}
		\cl(6) = \bigoplus_{k=0}^{6} \cl^k(6), \qquad
		\dim \cl^k(6) = \binom{6}{k},
		\label{eq:grade-decomp}
	\end{equation}
	has total dimension 64.  Of the seven grade subspaces, the relevant ones for
	this work are the bivectors $\cl^2(6)$ (dimension 15), which contain the Lie
	algebra of $\spin(6)\cong\su(4)$, and the pseudoscalar $\cl^6(6)$, spanned
	by $\Isix = e_1 e_2 e_3 e_4 e_5 e_6$ with $\Isix^2 = -1$.  The 8-dimensional
	minimal left ideal decomposes under $\su(3)$ as
	\begin{equation}
		\mathbf{8} = \mathbf{3} \oplus \bar{\mathbf{3}} \oplus \mathbf{1}_a
		\oplus \mathbf{1}_b,
		\label{eq:spinor-decomp}
	\end{equation}
	where $\mathbf{3}$ and $\bar{\mathbf{3}}$ are the color triplet and
	anti-triplet and $\mathbf{1}_a$, $\mathbf{1}_b$ are two color singlets.
	The full quark space is $\mathbb{C}^{24} = \mathbb{C}^8\otimes\mathbb{C}^3$,
	with $\mathbb{C}^8$ the $\cl(6)$ spinor space and $\mathbb{C}^3$ the
	standard QCD color space.
	
	\subsection{\texorpdfstring{$\su(3)$}{SU(3)} Embedding via
		\texorpdfstring{$\spin(6)\cong\su(4)$}{Spin(6)≅SU(4)}}
	\label{sec:embedding}
	
	The isomorphism chain
	\begin{equation}
		\su(3)_{\mathrm{color}} \subset \su(4) \cong \spin(6)
		\subset \cl^+(6) \subset \cl(6)
		\label{eq:embedding-chain}
	\end{equation}
	embeds $\su(3)$ as the stabilizer subgroup of a preferred spinor
	$\xi_0=(0,0,0,1)^T\in\mathbb{C}^4$.  The 15 bivectors
	$B_{ij}=(i/4)[e_i,e_j]$ span the Lie algebra of $\spin(6)$ and contain all
	eight $\su(3)$ generators as a subalgebra.  Gell-Mann normalization
	$\tr(T_a T_b)=\half\delta_{ab}$ is maintained; all 28 $\su(3)$ commutators
	are verified to maximum deviation $1.24\times10^{-16}$ (see
	Paper~1~\cite{FirestoneP1} Appendix~B).  Independently, the bivector
	$\su(3)$ octet---derived as the stabilizer of the $H=-\tfrac{3}{2}$
	eigenstate $|\xi_0\rangle$ within the 15-dimensional bivector
	space---is verified to close as a Lie algebra under the matrix
	commutator directly in the $8\times8$ representation (maximum
	closure residual $6.2\times10^{-16}$; Appendix~\ref{app:B}).
	$\su(3)$ gauge links act as
	standard $3\times3$ unitary matrices on the $\mathbb{C}^3$ color factor and
	are unchanged by the $\cl(6)$ embedding.  \textbf{$\su(3)$ is embedded by
		construction, following the known isomorphism chain; it is not derived from
		$\cl(6)$.}
	
	\subsection{The Cartan Hypercharge \texorpdfstring{$H$}{H}
		and Its Eigenspectrum}
	\label{sec:cartan}
	
	This section introduces algebraic content not present in
	Paper~1~\cite{FirestoneP1}.  The Cartan hypercharge of $\mathrm{spin}(6)$
	corresponding to the $\su(3)\times\mathrm{U}(1)$ Levi decomposition is
	\begin{equation}
		H = B_{12} + B_{34} + B_{56},
		\label{eq:cartan-H}
	\end{equation}
	where $B_{ij}=(i/4)[e_i,e_j]$ are bivectors.  In the tensor product
	representation~\eqref{eq:basis}, $H$ is diagonal:
	\begin{align}
		H &= -\half(\sigma_z\otimes I_2\otimes I_2)
		- \half(I_2\otimes\sigma_z\otimes I_2)
		- \half(I_2\otimes I_2\otimes\sigma_z)
		\label{eq:H-tensor} \\
		&= \half\,\mathrm{diag}(-3,-1,-1,+1,-1,+1,+1,+3).
		\label{eq:H-diag}
	\end{align}
	The eigenspectrum is $\{-\tfrac{3}{2},-\tfrac{1}{2},+\tfrac{1}{2},+\tfrac{3}{2}\}$
	with multiplicities $\{1,3,3,1\}$, matching the $\su(3)$ representation
	dimensions in~\eqref{eq:spinor-decomp}.  Lagrange spectral projectors onto
	each $H$-eigenspace are
	\begin{equation}
		P_r = \prod_{h' \neq h_r} \frac{H - h'\,I_8}{h_r - h'},
		\label{eq:lagrange-projector}
	\end{equation}
	where $h_r$ is the $H$-eigenvalue of representation $r$.  The resulting
	projectors satisfy $\tr(P_{\mathbf{3}})=3$, $\tr(P_{\bar{\mathbf{3}}})=3$,
	$\tr(P_{\mathbf{1}_a})=1$, $\tr(P_{\mathbf{1}_b})=1$, and
	$\sum_r P_r = I_8$, all to machine precision.  Idempotency
	$P_r^2 = P_r$ and orthogonality $P_r P_s = 0$ for $r\neq s$ also hold to
	machine precision.  Full verification in Appendix~\ref{app:C}.
	
	\subsection{The \texorpdfstring{$(\GaSeven, H)$}{(Γ₇, H)} Double Grading}
	\label{sec:double-grading}
	
	The eight spinor basis states are simultaneously labeled by two commuting
	operators.  The $\spin(6)$ chirality element
	\begin{equation}
		\GaSeven = i\Isix = i\cdot e_1 e_2 e_3 e_4 e_5 e_6
		= \sigma_z\otimes\sigma_z\otimes\sigma_z
		\label{eq:Gamma7}
	\end{equation}
	takes eigenvalues $\pm1$.  This is distinct from the 4D chiral operator
	$\gamma_5 = e_1 e_2 e_3 e_4 = -\sigma_z\otimes\sigma_z\otimes I_2$ used in
	the pion correlator.  The commutativity $[\GaSeven, H]=0$ holds identically
	in the representation~\eqref{eq:basis}, and the four $\su(3)$ representations
	are uniquely labeled by the pair $(\GaSeven, H)$:
	
	\begin{table}[ht]
		\centering
		\caption{The $(\GaSeven, H)$ double grading of the spinor basis states.
			This is the standard $\su(4)\supset\su(3)\times\mathrm{U}(1)$ Levi
			decomposition; see Ref.~\cite{Slansky1981}, Table~36.}
		\label{tab:double-grading}
		\begin{tabular}{@{}ccccp{4.5cm}@{}}
			\toprule
			\textbf{Rep} & $\GaSeven$ & $H$ & \textbf{dim} & \textbf{Basis states} \\
			\midrule
			$\mathbf{1}_a$ & $+1$ & $-3/2$ & 1
			& $|\!\uparrow\uparrow\uparrow\rangle$ \\[3pt]
			$\mathbf{3}$ & $+1$ & $+1/2$ & 3
			& $|\!\uparrow\downarrow\downarrow\rangle$,\;
			$|\!\downarrow\uparrow\downarrow\rangle$,\;
			$|\!\downarrow\downarrow\uparrow\rangle$ \\[3pt]
			$\bar{\mathbf{3}}$ & $-1$ & $-1/2$ & 3
			& $|\!\uparrow\uparrow\downarrow\rangle$,\;
			$|\!\uparrow\downarrow\uparrow\rangle$,\;
			$|\!\downarrow\uparrow\uparrow\rangle$ \\[3pt]
			$\mathbf{1}_b$ & $-1$ & $+3/2$ & 1
			& $|\!\downarrow\downarrow\downarrow\rangle$ \\
			\bottomrule
		\end{tabular}
	\end{table}
	
	The $(\GaSeven, H)$ labeling is the standard $\su(4)\supset\su(3)\times\mathrm{U}(1)$
	Levi decomposition, tabulated by Slansky~\cite{Slansky1981} (Table~36).
	With the normalization $Q_{B-L}=(2/3)H$, the triplets carry
	baryon-minus-lepton charge $Q_{B-L}=\pm1/3$ and the singlets carry
	$Q_{B-L}=\mp1$, matching the Pati--Salam $\su(4)_{\mathrm{color}}$
	model~\cite{PatiSalam1974}.  The same decomposition arises in Furey's
	algebraic classification
	program~\cite{Furey2012,Furey2014,Furey2015,Furey2018a,Furey2018b} as the
	identification of $\su(3)$ quantum numbers for a single generation of
	Standard Model fermions from the $\cl(6)$ minimal left ideal.
	
	The key distinction for the dynamics of this paper is that $\GaSeven\neq\gamma_5$.
	The relationship $\GaSeven = \gamma_5\cdot(ie_5 e_6)$ shows that $\GaSeven$
	combines the 4D Lorentz chirality with the chirality of the internal
	$(e_5,e_6)$ pair.  Both $\GaSeven$ and $\gamma_5$ commute with $H$ and with
	all four projectors $P_r$, but the triplet subspace $\mathbf{3}$ contains
	states of both $\gamma_5$ eigenvalues---a fact that is fundamental to the
	analysis of Sec.~\ref{sec:cl6-specific}.
	
	% ============================================================
	\section{Lattice Implementation}
	\label{sec:lattice}
	% ============================================================
	
	\subsection{Gauge Ensembles and Parameters}
	\label{sec:ensembles}
	
	All production calculations use hypercubic lattices with periodic boundary
	conditions in all four directions.  The production volume for the grade
	decomposition calculations is $N=6$ ($V=1296$ sites).  Two gauge couplings
	are employed: $\beta=4.5$ (confined phase, $\langle P\rangle\approx0.340$)
	and $\beta=6.0$ (deconfined phase, $\langle P\rangle\approx0.593$).  The
	phase boundary at $\beta_c\approx5.69$ was located in
	Paper~1~\cite{FirestoneP1}~\S4.1.2 and is not remeasured here.
	
	Gauge configurations are generated by Metropolis updates at acceptance
	rate $\approx50\%$.  Thermalization was extended to $\Ntherm=1000$ sweeps
	at $\beta=4.5$ and $\Ntherm=800$ sweeps at $\beta=6.0$ (compared to
	$\Ntherm=300$ sweeps in Paper~1~\cite{FirestoneP1}); this additional
	thermalization was adopted following the checkerboard correction described
	in Sec.~\ref{sec:checkerboard}.  Decorrelation uses $N_{\mathrm{decorr}}=5$
	sweeps.  Production ensembles: $\Ncfg=20$ configurations per $\beta$ value.
	
	Table~\ref{tab:plaquette-gate} reports the plaquette equilibration gate;
	both ensembles pass the $\pm1\%$ gate at $1\sigma$.
	
	\begin{table}[ht]
		\centering
		\caption{Plaquette equilibration gate.  Corrected even-odd sweep produces
			$\langle P\rangle$ within $1\%$ of Paper~1~\cite{FirestoneP1} $N=6$
			values in both phases.}
		\label{tab:plaquette-gate}
		\begin{tabular}{@{}ccccc@{}}
			\toprule
			$\beta$ & $\langle P\rangle$ measured
			& Paper~1 $N=6$ $\langle P\rangle$
			& Deviation & Gate \\
			\midrule
			4.5 & $0.34063\pm0.00050$ & $0.33816$ & $+0.73\%$
			& $\checkmark$~\textsc{pass} \\
			6.0 & $0.59580\pm0.00073$ & $0.59334$ & $+0.41\%$
			& $\checkmark$~\textsc{pass} \\
			\bottomrule
		\end{tabular}
	\end{table}
	
	\subsection{Wilson--Dirac Operator}
	\label{sec:wilson-dirac}
	
	The Wilson--Dirac operator in $\cl(6)$ is unchanged from
	Paper~1~\cite{FirestoneP1}~\S3.3:
	\begin{equation}
		\begin{split}
			(\dW\Psi)(x) = (m_0+4r)\,\Psi(x)
			&- \tfrac{1}{2}\sum_\mu\bigl[(1-re_\mu)\,U_\mu(x)\,\Psi(x+\hat{\mu}) \\
			&\quad + (1+re_\mu)\,U_\mu^\dagger(x-\hat{\mu})\,\Psi(x-\hat{\mu})\bigr],
		\end{split}
		\label{eq:wilson-dirac}
	\end{equation}
	with $r=1$ fixed and the action on $\mathbb{C}^8\otimes\mathbb{C}^3$ per site
	given by $(1\pm e_\mu)\otimes U_\mu(x)$.  The bare quark mass parameter for
	the heavy-mass production runs is $m_0=0.05$, corresponding to
	$\kappa=1/(2(m_0+4))=0.12346$.
	
	The conjugate-gradient solver is applied to $\dW^\dagger\dW\,x = \dW^\dagger b$
	with convergence criterion $\|\mathrm{r}\|/\|b\|<10^{-8}$ and maximum 500
	iterations.  Mean CG iterations were 18.95 at $\beta=4.5$ and 21.85 at
	$\beta=6.0$; there were zero CG failures in the $\Ncfg=20$ production
	ensembles at $m_0=0.05$.  The pion propagator is built from a single point-to-all solve at the
	origin, spanning all 24 spin--color source components ($8$ spinor $\times$
	$3$ color); the solve is exact rather than stochastic, so each configuration
	yields one deterministic correlator.  The three classes of measurement in
	this work accordingly use distinct source types: the Experiment~A grade
	fractions use this exact point-to-all solve; the $\kcren$/PCAC determination
	uses $Z_2$ stochastic sources; and the free-field and completeness
	diagnostics use a single $Z_2$ source, valid there because those identities
	are source-independent.
	
	\subsection{Even-Odd Checkerboard Sweep Correction}
	\label{sec:checkerboard}
	
	A methodological correction relative to some intermediate runs is documented
	here.  Earlier runs employed a vectorized parallel Metropolis sweep that
	computed all staples simultaneously before proposing any updates.  This
	approach violates detailed balance: once a proposal at site $x$ is accepted,
	its neighbors' staples are modified, but those neighbors were already
	evaluated against pre-update values within the same sweep.  The resulting
	bias scales with coupling strength and manifests as a systematic plaquette
	deficit---approximately 5.1\% at $\beta=4.5$ and 2.4\% at $\beta=6.0$,
	independent of sweep count.
	
	The corrected implementation uses an even-odd checkerboard sweep (see
	Paper~1~\cite{FirestoneP1}~\S3.4 for the Metropolis update framework).  For
	each direction $\mu$ and parity $p\in\{0,1\}$: staples for all parity-$p$
	sites are computed using the current link configuration (including
	freshly-updated parity-$(1-p)$ sites from the preceding half-step), and
	proposals and acceptances for parity-$p$ sites are made in parallel.  Within
	a parity class, no two sites share a same-direction, same-parity staple, so
	detailed balance holds exactly at each half-step.  All production data in
	this paper use the corrected checkerboard sweep.  We note that this
	correction retroactively suggests that some intermediate-$\beta$ plaquette
	values reported in Paper~1~\cite{FirestoneP1} ($\beta=5.0$ at 7\% low,
	$\beta=5.5$ at 5\% low) may have the same origin and warrant re-examination.
	
	% ============================================================
	\section{Grade Projectors: Construction and Verification}
	\label{sec:projectors}
	% ============================================================
	
	\subsection{Construction as Lagrange Spectral Projectors}
	\label{sec:proj-construction}
	
	The four grade projectors are constructed as Lagrange spectral projectors of
	the Cartan hypercharge $H$ defined in~\eqref{eq:cartan-H}.  The four
	$H$-eigenvalues are $h_{\mathbf{1}_a}=-\tfrac{3}{2}$,
	$h_{\mathbf{3}}=+\tfrac{1}{2}$, $h_{\bar{\mathbf{3}}}=-\tfrac{1}{2}$, and
	$h_{\mathbf{1}_b}=+\tfrac{3}{2}$; the projector onto representation $r$ is
	given by~\eqref{eq:lagrange-projector}.
	
	
	This construction guarantees idempotency $P_r^2=P_r$ and orthogonality
	$P_r P_s=0$ for $r\neq s$ by Lagrange interpolation theory, independently of
	the specific form of $H$.  Completeness $\sum_r P_r = I_8$ follows from the
	Lagrange basis being a partition of unity over any finite set of distinct
	points.  Full algebraic verification is reported in
	Appendix~\ref{app:C}, Table~\ref{tab:verif-C1}.
	
	The charge conjugation matrix $C$ in $\cl(6)$ is constructed from first
	principles via the requirement $C e_i C^\dagger = -e_i^*$ for all
	$i\in\{1,\ldots,6\}$.  It is not the standard $4\times4$ Dirac charge
	conjugation matrix; it must be derived from the $8\times8$ $\cl(6)$
	basis~\eqref{eq:basis}.  The resulting $C$ satisfies
	$C P_{\mathbf{3}} C^\dagger = P_{\bar{\mathbf{3}}}$ and
	$C P_{\mathbf{1}_a} C^\dagger = P_{\mathbf{1}_b}$ exactly (residual
	$0.00\times10^0$); see Appendix~\ref{app:B}.
	
	\subsection{\texorpdfstring{$\su(3)$}{SU(3)} Covariance}
	\label{sec:proj-covariance}
	
	The projectors commute with all eight $\su(3)$ generators:
	$[T_a, P_r]=0$ for all $a\in\{1,\ldots,8\}$ and all $r$.  This is not a
	separately proved identity but follows automatically from the algebraic
	structure: $H$ is a polynomial in the Cartan generators of $\spin(6)$, and
	the $\su(3)$ generators $T_a$ lie in the subalgebra that commutes with $H$
	by definition of the Cartan decomposition.  Numerically, $\su(3)$ covariance
	$U P_r U^\dagger = P_r$ is verified over 50 random $U\in\su(3)$ for all
	four projectors, with maximum deviation $2.44\times10^{-15}$
	(Table~\ref{tab:verif-C1}, row~C5).
	
	\subsection{Commutativity with
		\texorpdfstring{$\gamma_5$}{gamma5}}
	\label{sec:proj-gamma5}
	
	The commutativity $[\gamma_5, P_r]=0$ for all $r$ holds exactly (residual
	$0.00\times10^0$).
	This property is essential for the grade-projected pion correlator to be
	well-defined: without it, the operator $P_r\gamma_5 P_r$ would mix
	chiralities and $C_\pi^{(r)}(t)$ would not define a physical pseudoscalar
	channel.  The commutativity follows from the observation that
	$\gamma_5 = e_1 e_2 e_3 e_4$ acts only on the first two tensor factors
	in~\eqref{eq:basis}, while $H$ acts on all three; the commutator vanishes
	because $\gamma_5$ commutes with each $B_{ij}$ appearing in $H$.
	
	\subsection{\texorpdfstring{$H$}{H}-Grade Labeling via
		\texorpdfstring{$\Isix$}{I6}}
	\label{sec:proj-I6}
	
	The pseudoscalar $\Isix = e_1 e_2 e_3 e_4 e_5 e_6$ acts on the two singlet
	subspaces as
	\begin{equation}
		\begin{aligned}
			\Isix\,P_{\mathbf{1}_a} &= -i\,P_{\mathbf{1}_a}, \\
			\Isix\,P_{\mathbf{1}_b} &= +i\,P_{\mathbf{1}_b},
		\end{aligned}
		\label{eq:I6-singlet}
	\end{equation}
	verified to residual $0.00\times10^0$.
	This provides an alternative labeling of the singlet channels consistent with
	the $(\GaSeven, H)$ double grading of Sec.~\ref{sec:double-grading}.
	
	\subsection{Complete Verification Table}
	\label{sec:proj-verification}
	
	The full projector algebraic verification is reported as
	Table~\ref{tab:verif-C1} (Appendix~\ref{app:C}).  All five algebraic checks
	(C1--C5) and one stochastic check (S1) pass to machine precision.  The
	$I_4{\otimes}U(2)$ symmetry argument of Sec.~\ref{sec:kinematic-theorem}
	depends only on checks C1--C5; S1 is an additional stochastic verification
	of group-averaging behavior under the Haar measure.
	
	% ============================================================
	\section{Non-Perturbative \texorpdfstring{$\kcren$}{kappa_c^ren}
		Determination}
	\label{sec:kappa-c}
	% ============================================================
	
	\subsection{PCAC Method}
	\label{sec:pcac-method}
	
	The renormalized critical hopping parameter $\kcren(\beta)$ is the value of
	$\kappa$ at which the PCAC quark mass vanishes:
	\begin{equation}
		\mPCAC(\kappa) = \frac{C_{AP}(t)}{4\,C_{PP}(t)},
		\label{eq:mPCAC}
	\end{equation}
	where $C_{AP}(t)$ and $C_{PP}(t)$ are the axial-pseudoscalar and
	pseudoscalar-pseudoscalar correlators, respectively, plateaued over
	$t\in[N/4,\,3N/4]$.  In the $\cl(6)$ context, the axial current is
	$A_4^a(x,t) = \bar\psi(x)(e_4\otimes I_3)(\gamma_5\otimes I_3)
	(I_8\otimes\lambda^a/2)\psi(x)$, which requires the tensor product structure
	of Paper~1~\cite{FirestoneP1}~\S2.3.  The pseudoscalar density is
	$P^a(x,t) = \bar\psi(x)(\gamma_5\otimes I_3)(I_8\otimes\lambda^a/2)\psi(x)$.
	
	The zero crossing of $\mPCAC(\kappa)$ is located by Lagrange interpolation
	between the two $\kappa$ values bracketing $\mPCAC=0$, with the bracket
	selection requiring $\mathrm{SNR}>0.8$ on both endpoints to exclude
	noise-dominated crossings.  $Z_2$ stochastic sources are used for the
	correlator estimates.
	
	Method~B ($m_\pi^2$ extrapolation via cosh fit) was attempted but proved
	unreliable on $N\leq6$ lattices due to the insufficient plateau window at
	small volumes.  All $\kcren$ values quoted in this section use Method~A
	(PCAC) only.
	
	Because Method~B did not yield a usable determination, the $\kcren$ values
	reported here rest on the single PCAC method and carry no independent
	cross-check: the two-method consistency test that would normally validate a
	$\kcren$ determination could not be performed.  Accordingly these values are
	used only to locate the chiral regime for Experiment~B
	(Sec.~\ref{sec:experiment-B-design}); the $m_\pi^2$-extrapolation cross-check
	is identified as follow-up work in Sec.~\ref{sec:future-directions}.
	
	\subsection{Results}
	\label{sec:kappa-c-results}
	
	Block~1 of the computation produced partial results sufficient for the grade
	decomposition program.  Two of six $\beta$ values were determined; the
	remaining four are not required for Paper~2.
	
	\begin{table}[ht]
		\centering
		\caption{Non-perturbative $\kcren$ values used in this paper.
			The $\beta=4.5$ value is systematic-dominated; see text.}
		\label{tab:kappa-c}
		\begin{tabular}{@{}clp{7.5cm}@{}}
			\toprule
			$\beta$ & $\kcren$ & Source and notes \\
			\midrule
			4.5 & $0.1468\pm0.0050$
			& $N=6$, Method~A (PCAC); systematic-dominated.
			$N=4$/$N=6$ spread ($0.127\to0.147$) propagated as
			$\pm0.005$ systematic uncertainty. \\[4pt]
			6.0 & $0.1513\pm0.0006$
			& Weighted mean across four independent runs;
			$0.9\sigma$ from literature value
			$0.1518\pm0.0005$
			(Paper~1~\cite{FirestoneP1}~\S5.5.1). \\
			\bottomrule
		\end{tabular}
	\end{table}
	
	Monotonicity check: $\kcren(\beta=4.5)=0.1468 < 0.1513=\kcren(\beta=6.0)$,
	consistent with the physical expectation that weaker coupling (larger $\beta$)
	produces smaller additive mass renormalization and $\kcren$ closer to the
	free-field value $0.125$.  The large uncertainty on $\kcren(\beta=4.5)$
	reflects genuine $O(a)$ discretization effects at strong coupling that are
	not yet resolved; this uncertainty will propagate as a systematic into any
	chiral-regime calculation using this value.
	
	The present heavy-mass production runs at $m_0=0.05$ ($\kappa=0.12346$)
	lie well below $\kcren$ at both $\beta$ values and do not require the
	chiral-regime solver mitigations of Appendix~\ref{app:D}.  The
	$\kcren$ values are provided as prerequisites for Experiment~B
	(Sec.~\ref{sec:experiment-B-design}) and are referenced in
	Sec.~\ref{sec:outlook}.
	
	% ============================================================
	\section{Grade Decomposition of the Pion Correlator}
	\label{sec:grade-decomp}
	% ============================================================
	
	\subsection{The Grade-Projected Pion Correlator}
	\label{sec:correlator-def}
	
	The standard pion correlator in the $\cl(6)$ framework is
	\begin{equation}
		C_\pi(t) = \sum_{\mathbf{x}}\Bigl\langle
		\tr\bigl[S(\mathbf{x},t;\,\mathbf{0},0)\,\gamma_5\,
		S^\dagger(\mathbf{x},t;\,\mathbf{0},0)\,\gamma_5\bigr]
		\Bigr\rangle,
		\label{eq:pion-correlator}
	\end{equation}
	where $S(\mathbf{x},t;\mathbf{0},0) = \dW^{-1}(\mathbf{x},t;\mathbf{0},0)$
	is the quark propagator (a $24\times24$ matrix per site pair in
	$\mathbb{C}^8\otimes\mathbb{C}^3$), the trace is over both spinor and color
	indices, and $\gamma_5$ acts on the spinor sector only.  The grade-projected
	pion correlator for representation $r$ is defined by inserting the projector
	$P_r$ into the spinor trace:
	\begin{multline}
		C_\pi^{(r)}(t) = \sum_{\mathbf{x}}\Bigl\langle
		\tr\bigl[(P_r\otimes I_3)\,\gamma_5\,(P_r\otimes I_3) \\
		\cdot\,S(\mathbf{x},t;\,\mathbf{0},0)\,\gamma_5\,
		S^\dagger(\mathbf{x},t;\,\mathbf{0},0)\bigr]
		\Bigr\rangle.
		\label{eq:grade-projected-correlator}
	\end{multline}
	The commutativity $[\gamma_5, P_r]=0$ of Sec.~\ref{sec:proj-gamma5} ensures
	that $P_r$ and $\gamma_5$ can be commuted freely, so $C_\pi^{(r)}(t)$ defines
	a valid pseudoscalar channel.  The grade fraction is
	\begin{equation}
		f_r(t) \;\equiv\; \frac{C_\pi^{(r)}(t)}{C_\pi(t)}.
		\label{eq:grade-fraction}
	\end{equation}
	By the completeness identity $\sum_r P_r = I_8$, it follows algebraically that
	\begin{equation}
		\sum_r C_\pi^{(r)}(t) = C_\pi(t)
		\quad\text{(exact, algebraically)},
		\label{eq:correlator-completeness}
	\end{equation}
	and therefore $\sum_r f_r(t) = 1$ on every configuration and at every
	timeslice.  This identity provides a per-configuration sanity check: any
	violation indicates a bug in the projector construction or trace evaluation.
	
	\subsection{Algebraic Completeness Verification}
	\label{sec:completeness-verification}
	
	The identity~\eqref{eq:correlator-completeness} is verified via the
	per-configuration relative residual
	\begin{equation}
		\delta(t) = \frac{\bigl|C_\pi(t) - \sum_r C_\pi^{(r)}(t)\bigr|}
		{\bigl|C_\pi(t)\bigr|}.
		\label{eq:completeness-residual}
	\end{equation}
	In the free-field case ($U_\mu=I_3$ everywhere, $N=4$, $m_0=0.05$, single
	$Z_2$ source), $\delta_{\max} = 1.26\times10^{-16}$ over
	$t\in\{0,1,2,3\}$---six orders of magnitude below the target threshold
	of $10^{-10}$.  In the production ensembles ($N=6$, $\Ncfg=20$, a single
	point-to-all solve per configuration spanning 24 spin--color components):
	\begin{align*}
		\delta_{\max}(\beta=4.5)
		&= 5.371\times10^{-16}
		&&\text{(max over 20 configs $\times$ 6 timeslices),}\\
		\delta_{\max}(\beta=6.0)
		&= 5.366\times10^{-16}
		&&\text{(max over 20 configs $\times$ 6 timeslices).}
	\end{align*}
	Both values are consistent with floating-point roundoff at
	$O(\varepsilon_{\mathrm{mach}}\times N_{\mathrm{hops}})$.  The
	identity~\eqref{eq:correlator-completeness} is verified to machine precision
	on every configuration in the production ensemble.
	
	\subsection{Charge Conjugation and \texorpdfstring{$\gamma_5$}{gamma5}
		Identities}
	\label{sec:additional-idents}
	
	Three additional algebraic identities are verified to machine precision prior
	to the production runs (Appendix~\ref{app:B}):
	
	\begin{itemize}
		\item Charge conjugation pairing: $CP_3 C^\dagger = P_{\bar{3}}$ and
		$CP_{1_a} C^\dagger = P_{1_b}$, each to residual $0.00\times10^0$.
		\item $\gamma_5$ commutativity: $[\gamma_5, P_r] = 0$ for all $r$, to residual
		$0.00\times10^0$.
		\item $\Isix$ action on singlets: $\Isix\,P_{1_a} = -i\,P_{1_a}$ and
		$\Isix\,P_{1_b} = +i\,P_{1_b}$, each to residual $0.00\times10^0$.
	\end{itemize}
	
	These identities, together with the projector checks of
	Appendix~\ref{app:C}, establish the algebraic infrastructure to machine
	precision before any gauge configurations are introduced.
	
	\subsection{\texorpdfstring{$H$}{H} Eigenvalue and Additional Verifications}
	\label{sec:H-eigenvalue-verif}
	
	The Cartan hypercharge $H$ is verified to have eigenvalues
	$\{-\tfrac{3}{2},\,-\tfrac{1}{2},\,+\tfrac{1}{2},\,+\tfrac{3}{2}\}$ with
	multiplicities $\{1,3,3,1\}$ in the explicit basis~\eqref{eq:basis}.
	The commutativity $[\GaSeven,\,H]=0$ holds identically.  The commutativity
	of $H$ with all eight $\su(3)$ generators,
	\[
	[H,\,T_a] = 0
	\quad\text{for } a\in\{1,\ldots,8\},
	\]
	is verified numerically to maximum deviation $<10^{-15}$.  These
	verifications confirm that $H$ generates the $\mathrm{U}(1)$ factor of the
	$\su(3)\times\mathrm{U}(1)$ Levi subgroup, as required by the
	representation-theoretic argument of Sec.~\ref{sec:double-grading}.
	
	% ============================================================
	%  REMAINING SECTIONS — TO BE ADDED SECTION BY SECTION
	% ============================================================
	
	\subsection{The Kinematic Limit Theorem}
	\label{sec:kinematic-theorem}
	
	The central theoretical result of this paper is the following theorem,
	which governs the free-field behavior of the grade fractions and provides
	the null hypothesis for Experiment~B (Sec.~\ref{sec:experiment-B-design}).
	
	\begin{theorem}[Kinematic Limit]
		\label{thm:kinematic}
		Let $G_{\mathrm{free}} = (\dW^{\mathrm{free}})^{-1}$ be the free-field
		quark propagator (all gauge links $U_\mu = I_3$).  Then
		\begin{equation}
			f_r^{\mathrm{free}}(t)
			\;=\; \frac{C_\pi^{(r)}{}_{\!\mathrm{free}}(t)}{C_\pi^{\mathrm{free}}(t)}
			\;=\; \frac{\dim(r)}{8}
			\label{eq:kinematic-free}
		\end{equation}
		exactly, for all $r\in\{3,\,\bar{3},\,1_a,\,1_b\}$ and all timeslices $t$.
		In the interacting theory, the same identity holds in the ensemble limit
		\begin{equation}
			\lim_{\Ncfg\to\infty}\langle f_r(t)\rangle \;=\; \frac{\dim(r)}{8},
			\label{eq:kinematic-interacting}
		\end{equation}
		enforced by $\su(3)$ color symmetry and the hypercubic symmetry of the
		ensemble measure.
	\end{theorem}
	
	\begin{proof}
		The proof requires two ingredients.  The first is the commutativity
		$[\gamma_5, P_r]=0$ established in Sec.~\ref{sec:proj-gamma5}: without it,
		the grade-projected pion operator $P_r\gamma_5 P_r$ would not define a valid
		pseudoscalar channel.  The second concerns the Hermitian matrix
		\begin{equation}
			\Qtil(t) = \sum_{\mathbf{x}}(\gamma_5\otimes I_3)\,
			G(\mathbf{x},t;\mathbf{0},0)\,(\gamma_5\otimes I_3)\,
			G^\dagger(\mathbf{x},t;\mathbf{0},0),
			\label{eq:Qtil-def}
		\end{equation}
		whose grade-projected traces give the channel correlators,
		$C_\pi^{(r)} = \tr[(P_r\otimes I_3)\Qtil]$ and $C_\pi = \tr[\Qtil]$.  Since
		each $P_r$ is diagonal in the $(\GaSeven,H)$ grade basis (in which $H$ is
		diagonal), these traces depend only on the diagonal of $\Qtil$ in that
		basis.  We show that for the free field this diagonal is constant:
		\begin{equation}
			\bigl[\Qtil_{\mathrm{free}}(t)\bigr]_{\text{grade diag}} = c(t)\,I_{24}.
			\label{eq:Qtil-diag}
		\end{equation}
		Given~\eqref{eq:Qtil-diag},
		\[
		C_\pi^{(r)} = c\,\tr[P_r\otimes I_3] = 3c\,\dim(r),
		\qquad
		C_\pi = 24c,
		\]
		and $f_r = \dim(r)/8$ follows immediately, for all $r$ and all $t$.
		
		It remains to establish~\eqref{eq:Qtil-diag}.  In the tensor-product basis
		of Sec.~\ref{sec:basis} the spacetime generators satisfy
		$e_\mu = M_\mu\otimes I_2$ for all $\mu\in\{1,2,3,4\}$: the third Clifford
		qubit factor is always $I_2$.  Hence $\dW^{\mathrm{free}}$ commutes with
		every $I_4\otimes U''\otimes I_3$ for $U''\in\mathrm{U}(2)$, and by Schur's
		lemma ($\mathrm{U}(2)$ acts irreducibly on the third $\mathbb{C}^2$) any
		matrix commuting with all such elements has the form $X\otimes I_2$.  In the
		free theory the propagator is color-diagonal ($U_\mu=I_3$), so
		$X = A_4\otimes I_3$, with $A_4$ acting on the four-dimensional $(q_1,q_2)$
		subspace.  Summing over spatial sinks at fixed $t$ projects onto zero spatial
		momentum, where the residual invariance is the \emph{spatial-cubic group
		$\times$ time reflection}, not the full hypercubic group.  Spatial-cubic
		invariance removes the spatial generators $e_2,e_3,e_4$ (which form a cubic
		vector and average to zero) together with all tensor structures, leaving
		\begin{equation}
			A_4(t) = \alpha(t)\,I_4 + \beta(t)\,e_1,
			\label{eq:A4-form}
		\end{equation}
		where $e_1=\sigma_x\otimes I_2$ is the time-direction generator.  The term
		$\beta\,e_1$ need not vanish---time reflection forces $\beta(t)$ to be odd
		in $t$, so it vanishes only at the fixed timeslices $t=0$ and $t=L/2$---but
		$e_1$ has \emph{identically zero diagonal} in the grade basis and so
		contributes nothing to the diagonal of $\Qtil = A_4\otimes I_2\otimes I_3$,
		which is therefore $\alpha(t)\,I_{24}$.  This
		establishes~\eqref{eq:Qtil-diag}; the grade-fraction result rests only on
		this zero-diagonal property and holds whether or not $\beta=0$.
		
		The interacting identity~\eqref{eq:kinematic-interacting} is a distinct,
		ensemble-averaged statement: it follows from applying the same $\su(3)$
		color and spatial-cubic symmetries to $\langle\Qtil\rangle$, and holds for
		$\langle f_r\rangle$ in the $\Ncfg\to\infty$ limit rather than
		configuration-by-configuration.
	\end{proof}	
	The physical content of the theorem is that the free Wilson--Dirac operator
	generates no dynamics in the $e_5, e_6$ directions: those generators appear
	nowhere in $\dW^{\mathrm{free}}$.  The grade decomposition is defined by
	$H = B_{12}+B_{34}+B_{56}$, and the bivector
	$B_{56} = \tfrac{i}{4}[e_5,e_6]$ is algebraically orthogonal to everything
	$\dW^{\mathrm{free}}$ can do.
	
	\medskip
	\noindent\textit{Numerical verification.}\enspace
	The theorem is verified at machine precision on the free lattice
	($U_\mu = I_3$ everywhere, $N=4$, $m_0=0.05$, single $Z_2$ source):
	\begin{center}
		\begin{tabular}{@{}lccc@{}}
			\toprule
			Grade & $f_r$ measured & $\dim(r)/8$ expected & Residual \\
			\midrule
			$\mathbf{3}$         & $0.375000$ & $0.375000$ & $8.5\times10^{-15}$ \\
			$\bar{\mathbf{3}}$   & $0.375000$ & $0.375000$ & $8.5\times10^{-15}$ \\
			$\mathbf{1}_a$       & $0.125000$ & $0.125000$ & $8.4\times10^{-15}$ \\
			$\mathbf{1}_b$       & $0.125000$ & $0.125000$ & $8.4\times10^{-15}$ \\
			\bottomrule
		\end{tabular}
	\end{center}
	Maximum residual from the expected free-field values:
	$O(\varepsilon_{\mathrm{mach}})$ at all timeslices.
	
	
	% ============================================================
	%  REMAINING SECTIONS — TO BE ADDED SECTION BY SECTION
	% ============================================================
	
	\subsection{Experiment~A: Heavy-Mass Production Results}
	\label{sec:experiment-A}
	
	Experiment~A measures the grade fractions on thermalized quenched production
	ensembles at $N=6$, $m_0=0.05$, $\beta\in\{4.5,\,6.0\}$, with $\Ncfg=20$
	configurations per phase.  Parameters are as in Sec.~\ref{sec:ensembles};
	the bare quark mass $m_0=0.05$ corresponds to $\kappa=0.123457$, placing
	the quarks well below $\kcren$ at both $\beta$ values and deep in the
	heavy-mass regime ($m_{\mathrm{phys}}\approx0.76$ in lattice units).
	
	\begin{table}[ht]
		\centering
		\caption{Grade fractions $f_r$ at $t=1$.  Free-field reference:
			$f_r=\dim(r)/8\in\{0.375,\,0.375,\,0.125,\,0.125\}$.
			$\Delta\equiv f_r(\beta{=}4.5)-f_r(\beta{=}6.0)$.
			Errors are statistical, from a jackknife over the $\Ncfg=20$ gauge configurations (one point-to-all solve per configuration).}
		\label{tab:grade-fractions}
		\begin{tabular}{@{}llcccc@{}}
			\toprule
			Grade & Rep              & $f_r(\beta=4.5)$        & $f_r(\beta=6.0)$
			& $\Delta$                & $\sigma$ \\
			\midrule
			$\mathbf{3}$
			& Triplet
			& $0.374918\pm0.000079$
			& $0.374968\pm0.000028$
			& $-0.000049\pm0.000084$
			& $0.6$ \\[3pt]
			$\bar{\mathbf{3}}$
			& Anti-triplet
			& $0.375081\pm0.000080$
			& $0.375032\pm0.000028$
			& $+0.000049\pm0.000085$
			& $0.6$ \\[3pt]
			$\mathbf{1}_a$
			& Singlet $(-i)$
			& $0.125082\pm0.000079$
			& $0.125032\pm0.000028$
			& $+0.000049\pm0.000084$
			& $0.6$ \\[3pt]
			$\mathbf{1}_b$
			& Singlet $(+i)$
			& $0.124919\pm0.000080$
			& $0.124968\pm0.000028$
			& $-0.000049\pm0.000085$
			& $0.6$ \\
			\midrule
			Sum & Total
			& $1.000000$
			& $1.000000$
			& $0.000000$
			& — \\
			\bottomrule
		\end{tabular}
	\end{table}
	
	Maximum deviation from the free-field reference at $t=1$:
	$\approx8.2\times10^{-5}$
	at $\beta=4.5$.  All four phase differences are consistent with zero at
	$0.6\sigma$ significance.
	
	The charge-conjugation pairing is manifest in the data: $\Delta f_3 =
	-\Delta f_{\bar{3}}$ and $\Delta f_{1_a} = -\Delta f_{1_b}$ to within
	statistical precision, as required by the identity $CP_3 C^\dagger =
	P_{\bar{3}}$ established in Sec.~\ref{sec:additional-idents}.  The four
	$\sigma$-values group into two exactly-equal pairs:
	$\sigma(f_{1_a})=\sigma(f_3)$ and $\sigma(f_{1_b})=\sigma(f_{\bar 3})$ hold
	to machine precision, forced configuration-by-configuration by the exact
	identity of Proposition~\ref{prop:gamma7-frozen}
	($f_{1_a}=\half-f_3$, $f_{1_b}=\half-f_{\bar 3}$); the remaining equality
	$\sigma(f_3)\approx\sigma(f_{\bar 3})$ holds only statistically, by the
	ensemble-level charge-conjugation relation
	$\langle f_3\rangle=\langle f_{\bar 3}\rangle$.  Per configuration the four
	fractions therefore carry exactly two independent degrees of freedom,
	reducing to the single C-even observable $\ftriplet$ only in the ensemble
	average (Sec.~\ref{sec:cl6-specific}).
	
	\medskip
	\noindent\textit{$t$-dependence.}\enspace
	The grade fractions at all timeslices $t\in\{0,\ldots,5\}$ are provided in
	the extended Table~6.1 of the Data Tables.  Errors increase rapidly with $t$
	because the pion correlator decays as $e^{-m_\pi t}$ with $m_\pi\approx3.4$
	in lattice units: by $t=3$, the signal has fallen by $e^{-10}\approx
	5\times10^{-5}$.  An apparent $3.6\sigma$ deviation in $f_3$ at
	$\beta=6.0$, $t=3$ is noise-dominated with $\Ncfg=20$ and does not
	constitute a physical signal; this is documented further in
	Sec.~\ref{sec:cl6-specific}.  At $t=0$, the grade fractions satisfy
	$f_r\approx\dim(r)/8$ to within $\approx3\times10^{-6}$, consistent with
	the short-range nature of the propagator at small $t$.
	
	\subsection{Discussion}
	\label{sec:discussion}
	
	\subsubsection{The kinematic limit theorem}
	\label{sec:kinematic-discussion}
	
	The result $f_r\approx\dim(r)/8$ observed throughout Table~\ref{tab:grade-fractions}
	is not a counting argument.  It is a theorem with specific algebraic content
	(Theorem~\ref{thm:kinematic}, Sec.~\ref{sec:kinematic-theorem}), and
	identifying that content determines what the small deviations measure.
	
	The physical content is that $\dW^{\mathrm{free}}$ generates no dynamics in
	the $e_5, e_6$ directions: those generators appear nowhere in $\dW^{\mathrm{free}}$.
	The grade decomposition is defined by $H = B_{12}+B_{34}+B_{56}$, and the
	bivector $B_{56} = \tfrac{i}{4}[e_5,e_6]$ is algebraically orthogonal to
	everything $\dW^{\mathrm{free}}$ can do.  Uniform distribution across
	$H$-eigenspaces follows not from any symmetry between $e_5, e_6$ and the
	spacetime generators, but from the absence of any coupling between
	$\dW^{\mathrm{free}}$ and the internal directions that define the grading.
	
	\subsubsection{The \texorpdfstring{$(\GaSeven, H)$}{(Gamma7, H)} double
		grading and its place in the literature}
	\label{sec:grading-literature}
	
	Each grade projector is labeled by a pair of eigenvalues of two commuting
	operators.  The $\spin(6)$ chirality $\GaSeven = i\Isix$, where
	$\Isix = e_1 e_2 e_3 e_4 e_5 e_6$ is the $\cl(6)$ pseudoscalar, takes
	eigenvalues $\pm1$ and in the tensor product representation reads
	$\GaSeven = \sigma_z\otimes\sigma_z\otimes\sigma_z$.  The Cartan hypercharge
	$H = B_{12}+B_{34}+B_{56}$ takes eigenvalues
	$\{-\tfrac{3}{2},\,-\tfrac{1}{2},\,+\tfrac{1}{2},\,+\tfrac{3}{2}\}$ with
	multiplicities $\{1,3,3,1\}$.  Both operators are diagonal in the standard
	computational basis of $\mathbb{C}^8$, and their commutativity $[\GaSeven, H]=0$
	follows immediately.  The four representations are uniquely identified by the
	pair $(\GaSeven, H)$:
	\begin{center}
		\begin{tabular}{@{}ccccp{4.5cm}@{}}
			\toprule
			Rep & $\GaSeven$ & $H$ & dim & Basis states \\
			\midrule
			$\mathbf{1}_a$     & $+1$ & $-3/2$ & $1$
			& $|\!\uparrow\uparrow\uparrow\rangle$ \\[3pt]
			$\mathbf{3}$       & $+1$ & $+1/2$ & $3$
			& $|\!\uparrow\downarrow\downarrow\rangle$,
			$|\!\downarrow\uparrow\downarrow\rangle$,
			$|\!\downarrow\downarrow\uparrow\rangle$ \\[3pt]
			$\bar{\mathbf{3}}$ & $-1$ & $-1/2$ & $3$
			& $|\!\uparrow\uparrow\downarrow\rangle$,
			$|\!\uparrow\downarrow\uparrow\rangle$,
			$|\!\downarrow\uparrow\uparrow\rangle$ \\[3pt]
			$\mathbf{1}_b$     & $-1$ & $+3/2$ & $1$
			& $|\!\downarrow\downarrow\downarrow\rangle$ \\
			\bottomrule
		\end{tabular}
	\end{center}
	
	A distinction central to the interpretation is that $\GaSeven\neq\gamma_5$.
	The 4D chiral operator used in the pion correlator is
	$\gamma_5 = e_1 e_2 e_3 e_4 = -\sigma_z\otimes\sigma_z\otimes I_2$, while
	the $\spin(6)$ chirality is
	$\GaSeven = \sigma_z\otimes\sigma_z\otimes\sigma_z$.  Their relationship is
	$\GaSeven = \gamma_5\cdot(ie_5 e_6)$, reflecting the fact that $\GaSeven$
	combines the 4D Lorentz chirality with the chirality of the internal $(e_5,e_6)$
	pair.  Both commute with $H$ and with all four projectors $P_r$, but their
	action on the spinor space differs: the triplet and antitriplet subspaces are
	each invariant under $\gamma_5$ but are not $\gamma_5$ eigenspaces---a point
	returned to in Sec.~\ref{sec:cl6-specific}.
	
	The $(\GaSeven, H)$ double grading is a known structure in the representation
	theory of $\su(4)$.  The operator $H$ generates the $\mathrm{U}(1)$ factor of
	the Levi subgroup $\su(3)\times\mathrm{U}(1)\subset\su(4)$, where
	$\su(3) = \mathrm{Stab}_{\su(4)}(\xi_0)$ is the stabilizer of the preferred
	spinor direction $\xi_0=(0,0,0,1)^T$ that defines the color embedding.  With
	the normalization $Q_{B-L} = (2/3)H$, the triplet/antitriplet carry
	$Q_{B-L} = \pm1/3$ and the singlets carry $Q_{B-L} = \mp1$, matching the
	baryon-minus-lepton charges of the Pati--Salam $\su(4)_{\mathrm{color}}$
	model~\cite{PatiSalam1974}.  The complete $\su(4)\supset\su(3)\times\mathrm{U}(1)$
	branching table---$4\to3_{+1/3}\oplus1_{-1}$ and
	$\bar{4}\to\bar{3}_{-1/3}\oplus1_{+1}$---is tabulated in
	Ref.~\cite{Slansky1981} (Table~36).  In Furey's algebraic classification
	program~\cite{Furey2012,Furey2014,Furey2015,Furey2018a,Furey2018b}, the same
	decomposition identifies the $\su(3)$ quantum numbers of a single generation
	of Standard Model fermions from the $\cl(6)$ minimal left ideal.
	
	The contribution of the present work is the first dynamical implementation of
	this grading.  The projectors $P_r$ are constructed as Lagrange spectral
	projectors of $H$, verified to satisfy completeness, idempotency,
	orthogonality, $\su(3)$ covariance, and $\gamma_5$-commutativity to residual
	$\leq 2.4\times10^{-15}$ (Table~\ref{tab:verif-C1}), and used to decompose a
	physical pion correlator on a Monte Carlo gauge ensemble.  The algebraic
	completeness identity $\sum_r C_\pi^{(r)}(t) = C_\pi(t)$ holds to
	$\delta_{\max} = 5.4\times10^{-16}$ on every configuration.  We emphasize
	that the structural analogy between the singlet subspaces $\mathbf{1}_a,
	\mathbf{1}_b$ and lepton quantum numbers in the Furey program does not apply
	in the dynamical QCD context of this work: here, color is carried by the
	separate $\mathbb{C}^3$ factor of the full quark space
	$\mathbb{C}^{24} = \mathbb{C}^8\otimes\mathbb{C}^3$, and the grade projectors
	decompose the spinor sector $\mathbb{C}^8$ only.
	
	\subsubsection{Interpretation of the Experiment~A deviations}
	\label{sec:deviation-interp}
	
	The production data of Table~\ref{tab:grade-fractions} show grade fractions
	consistent with $\dim(r)/8$ at the $10^{-4}$ level.  All phase differences
	$\Delta f_r = f_r(\beta{=}4.5) - f_r(\beta{=}6.0)$ are consistent with zero
	at $0.6\sigma$ significance.  We identify three candidate mechanisms for the
	observed deviations from free-field values and assess their relative importance
	at the current parameter point.
	
	The primary driver is gauge-link color-spin entanglement.  The Wilson hopping
	term $(1\pm e_\mu)\otimes U_\mu(x)$ does not preserve $H$-eigenspaces: one
	finds $[H, e_\mu]\neq 0$ for all $\mu\in\{1,2,3,4\}$.  Explicitly,
	\begin{equation}
		[H,\, e_1]
		\;=\;
		\bigl[-\tfrac{1}{2}\sigma_z\otimes I_2\otimes I_2,\;
		\sigma_x\otimes I_2\otimes I_2\bigr]
		\;=\;
		-i\,\sigma_y\otimes I_2\otimes I_2
		\;\neq\; 0,
		\label{eq:H-e1-commutator}
	\end{equation}
	and analogously for $\mu=2,3,4$.  Since $P_r$ is a polynomial in $H$, it
	follows that $[P_r, \dW]\neq 0$: every lattice hop mixes grade channels, and
	the propagator $G = \dW^{-1}$ accumulates this mixing with each step.  On a
	single gauge configuration, $\Qtil_{\mathrm{int}}$ is a full $24\times24$
	matrix rather than a multiple of $I_{24}$, and the grade fractions can deviate
	from $\dim(r)/8$.  In the ensemble limit, $\su(3)$ and hypercubic symmetry
	restore $\langle\Qtil_{\mathrm{int}}\rangle\propto I_{24}$, but this
	convergence is slow: with $\Ncfg=20$, the residual per-configuration
	entanglement is incompletely averaged.
	
	Subleading contributions arise from the Wilson term's explicit $O(a)$ breaking
	of Lorentz symmetry, which introduces a small non-isotropic correction to
	$\Qtil_{\mathrm{int}}$ that is independent of $\Ncfg$ and vanishes only in the
	continuum limit.  Genuine non-perturbative vacuum dressing of grade
	channels---mechanism (c), the target of Experiment~B---is suppressed at
	$m_0=0.05$, where $m_{\mathrm{phys}}\approx0.76$ in lattice units and chiral
	effects scale as $(m_\pi/\Lambda_{\mathrm{QCD}})^2$.
	
	The observed sign pattern is consistent with the primary mechanism.  At $t=1$
	and $\beta=4.5$, the triplet fraction is depleted
	($\Delta f_3\approx-8.2\times10^{-5}$) while the singlet-$a$ fraction is
	correspondingly enhanced ($\Delta f_{1_a}\approx+8.2\times10^{-5}$).  This
	direction is expected from a perturbative $O(g^2)$ analysis: gauge-link
	dressing couples preferentially to color-charged spinor channels (the
	triplets), transferring a small fraction of weight toward the color-neutral
	singlet subspaces.
	
	The $t$-dependence carries additional diagnostic information but is currently
	noise-limited.  Deviations grow from $\sim10^{-4}$ at $t=1$ to
	$\sim7\times10^{-4}$ at $t=3$ at $\beta=6.0$, reflecting the accumulation of
	grade mixing over additional hops as expected from~\eqref{eq:H-e1-commutator}.
	However, at $m_0=0.05$ the pion correlator is suppressed as $e^{-m_\pi t}$
	with $m_\pi\approx3.4$ in lattice units, so $C_\pi(t{=}3)\sim
	e^{-10}\,C_\pi(0)$.  The apparent $3.6\sigma$ deviation at $t=3$ is
	signal-to-noise dominated with $\Ncfg=20$ and does not constitute a physical
	signal at current statistics; resolving the $t$-profile reliably requires
	$\Ncfg\gtrsim200$.
	
	\subsubsection{Channel-pair structure and the C-symmetry reduction}
	\label{sec:channel-pair}
	
	The algebraic charge-conjugation pairing $CP_3C^\dagger = P_{\bar{3}}$ and
	$CP_{1_a}C^\dagger = P_{1_b}$ (Appendix~\ref{app:B}), together with the
	completeness constraint $\sum_r f_r = 1$, determines the observable structure
	precisely.
	
	In the ensemble limit with the C-symmetric Wilson gauge action and a
	C-symmetric point source, charge conjugation enforces
	$\langle f_3\rangle = \langle f_{\bar{3}}\rangle$ and
	$\langle f_{1_a}\rangle = \langle f_{1_b}\rangle$.  Combined with completeness,
	this reduces four apparent observables to one independent C-even quantity:
	\begin{equation}
		\ftriplet \;\equiv\; f_3 + f_{\bar{3}},
		\label{eq:f-triplet}
	\end{equation}
	with $\fsinglet = 1 - \ftriplet$ determined by completeness.  The two C-odd
	combinations
	\begin{equation}
		A_{\mathrm{color}} = f_3 - f_{\bar{3}},
		\qquad
		A_{\mathrm{singlet}} = f_{1_a} - f_{1_b},
		\label{eq:C-odd-combos}
	\end{equation}
	both vanish in the $\Ncfg\to\infty$ ensemble limit for $\theta=0$.  Their
	observed magnitudes ($\sim1.6\times10^{-4}$ at $\Ncfg=20$) are finite-sample
	fluctuations expected to scale as $\Ncfg^{-1/2}$, and they serve as
	self-consistency checks rather than physics signals.  We note explicitly that
	$A_{\mathrm{color}}$ is C-odd and parity-even: it is not a probe of parity
	violation.  A nonzero ensemble-average $A_{\mathrm{color}}$ would require
	C violation; the $\theta$-term, being C-even, cannot generate it, and the
	present calculations are C-symmetric.  By
	Proposition~\ref{prop:gamma7-frozen} this is moot in any case, since
	$A_{\mathrm{color}} = -A_{\mathrm{singlet}}$ exactly per configuration.
	
	\begin{proposition}[$\spin(6)$ chirality is dynamically frozen]
	\label{prop:gamma7-frozen}
		On every gauge configuration and at every timeslice,
		\begin{equation}
			C_\pi^{(\GaSeven=+1)}(t) = C_\pi^{(\GaSeven=-1)}(t),
			\qquad\text{hence}\qquad
			f_{(\GaSeven=+1)}(t) = f_{(\GaSeven=-1)}(t) = \half ,
			\label{eq:gamma7-half}
		\end{equation}
		so that $A_{\GaSeven} \equiv A_{\mathrm{color}} + A_{\mathrm{singlet}} = 0$
		identically.  The result is independent of $\beta$, of $t$, and of the
		gauge field.
	\end{proposition}
	
	\begin{proof}
		The Wilson--Dirac operator is built only from the spacetime generators
		$e_1,\ldots,e_4$, each acting as $M_\mu\otimes I_2$ on the third Clifford
		qubit (Sec.~\ref{sec:basis}), while the gauge links multiply the color
		index alone.  The third qubit is therefore inert and the propagator
		factorizes as $G = \widehat{G}\otimes I_2$ on that qubit on \emph{every}
		configuration, free or interacting---a property unquenching cannot remove,
		as it alters the gauge measure but not this operator content.  The
		$\gamma_5$-dressed correlator bilinear $M(\mathbf{x},t) =
		(\gamma_5\otimes I_3)\,G\,(\gamma_5\otimes I_3)\,G^\dagger$ (the
		per-configuration summand of $\Qtil$) inherits the same factorization,
		$M = \widehat{M}\otimes I_2$, since $\gamma_5 = \widehat{\gamma}_5\otimes
		I_2$ likewise leaves the third qubit untouched.  With $\GaSeven =
		\sigma_z\otimes\sigma_z\otimes\sigma_z$, the grade-summed chiral asymmetry
		factorizes across the third qubit,
		\[
		C_\pi^{(\GaSeven=+1)} - C_\pi^{(\GaSeven=-1)}
		= \sum_{\mathbf{x}} \tr\!\bigl[(\GaSeven\otimes I_3)\,M\bigr]
		= \sum_{\mathbf{x}}
		\tr_{12c}\!\bigl[(\sigma_z\otimes\sigma_z\otimes I_3)\,\widehat{M}\bigr]\,
		\tr_{q_3}[\sigma_z] ,
		\]
		and the factor $\tr_{q_3}[\sigma_z]=0$ annihilates it on every
		configuration.  Dividing by $C_\pi$ and using the per-configuration
		completeness $\sum_r f_r = 1$ gives $f_{(\GaSeven=\pm1)} = \half$.
	\end{proof}
	
	This identity is structurally stronger than, and must not be conflated with,
	the charge-conjugation relation.  It fixes the \emph{sum} $A_{\mathrm{color}}
	+ A_{\mathrm{singlet}}$ to zero on \emph{every} configuration: the measured
	residual is $\max|A_{\GaSeven}| \le 3.9\times10^{-16}$ across both ensembles,
	all configurations and timeslices, i.e.\ $A_{\mathrm{color}} =
	-A_{\mathrm{singlet}}$ \emph{exactly}, not ``to within statistical
	precision.''  By contrast the individual asymmetries are each nonzero on a
	given configuration (per-configuration $|f_3 - f_{\bar 3}|$ has rms
	$\approx 1.1\times10^{-3}$); only their \emph{ensemble averages} vanish, by
	the genuinely statistical charge-conjugation relation $\langle f_3\rangle =
	\langle f_{\bar 3}\rangle$, $\langle f_{1_a}\rangle = \langle f_{1_b}\rangle$.
	The $\spin(6)$-chirality direction therefore carries no dynamical information
	whatsoever (Sec.~\ref{sec:cl6-specific}).
	
	The surviving observable $\ftriplet = f_3 + f_{\bar{3}}$ measures the fraction
	of pion correlator weight in the color-carrying spinor channels.  Its
	free-field value is $3/4$, verified in both phases at the $0.6\sigma$ level
	(Table~\ref{tab:grade-fractions}).  The phase difference
	$\dftriplet \equiv \ftriplet(\beta{=}4.5) - \ftriplet(\beta{=}6.0)$ is the
	target observable of Experiment~B: it measures whether the gauge vacuum
	distributes spinor weight differently between color-carrying and color-neutral
	internal channels at the chiral point, and constitutes the grade-decomposed
	analog of $\Delta A$ from Paper~1~\cite{FirestoneP1}.
	
	\subsubsection{The grade decomposition as a \texorpdfstring{$\cl(6)$}{Cl(6)}-specific
		structure}
	\label{sec:cl6-specific}
	
	A natural question is whether the grade-projected correlator $C_\pi^{(r)}(t)$
	could be computed in the standard four-component Wilson lattice QCD formulation
	by inserting an appropriate projector.  The answer is no, for a specific
	algebraic reason.
	
	The grade projectors are Lagrange spectral projectors of
	$H = B_{12}+B_{34}+B_{56}$, where
	\begin{equation}
		B_{56} \;=\; \frac{i}{4}[e_5,\, e_6].
		\label{eq:B56}
	\end{equation}
	The bivector $B_{56}$ requires the internal Clifford generators $e_5$ and
	$e_6$, which belong to $\cl(6)$ but have no counterpart in the standard
	four-component Dirac algebra $\cl(4)$.  The generators available in the
	standard formulation are $\{\gamma_1,\gamma_2,\gamma_3,\gamma_4\}$ and their
	products; without $e_5$ and $e_6$, the bivector $B_{56}$ cannot be formed,
	$H$ cannot be constructed, and therefore $P_r$ cannot be defined.
	
	A further obstruction arises from the structure of the $\gamma_5$ chirality
	decomposition.  The triplet subspace $\mathbf{3}$ (the $H=+1/2$ eigenspace)
	contains two basis states with $\gamma_5=+1$ and one with $\gamma_5=-1$; the
	antitriplet $\bar{\mathbf{3}}$ similarly has mixed $\gamma_5$ eigenvalues.
	The explicit assignments are:
	\begin{center}
		\begin{tabular}{@{}ccccc@{}}
			\toprule
			State & Rep & $\gamma_5$ & $\GaSeven$ & $H$ \\
			\midrule
			$|\!\uparrow\downarrow\downarrow\rangle$
			& $\mathbf{3}$ & $+1$ & $+1$ & $+1/2$ \\
			$|\!\downarrow\uparrow\downarrow\rangle$
			& $\mathbf{3}$ & $+1$ & $+1$ & $+1/2$ \\
			$|\!\downarrow\downarrow\uparrow\rangle$
			& $\mathbf{3}$ & $-1$ & $+1$ & $+1/2$ \\[3pt]
			$|\!\uparrow\uparrow\downarrow\rangle$
			& $\bar{\mathbf{3}}$ & $-1$ & $-1$ & $-1/2$ \\
			$|\!\uparrow\downarrow\uparrow\rangle$
			& $\bar{\mathbf{3}}$ & $+1$ & $-1$ & $-1/2$ \\
			$|\!\downarrow\uparrow\uparrow\rangle$
			& $\bar{\mathbf{3}}$ & $+1$ & $-1$ & $-1/2$ \\[3pt]
			$|\!\uparrow\uparrow\uparrow\rangle$
			& $\mathbf{1}_a$ & $-1$ & $+1$ & $-3/2$ \\
			$|\!\downarrow\downarrow\downarrow\rangle$
			& $\mathbf{1}_b$ & $-1$ & $-1$ & $+3/2$ \\
			\bottomrule
		\end{tabular}
	\end{center}
	No combination of standard Dirac projectors---chiral projectors
	$(1\pm\gamma_5)/2$, helicity projectors, or spin projectors
	$\sigma_{\mu\nu}$---can isolate a grade subspace, because the grade structure
	cuts across the $\gamma_5$ sectors.  This is not an artifact of the chosen
	basis; it is a consequence of the tensor product structure of the
	eight-dimensional spinor and cannot be removed by any change of basis within
	$\cl(4)$.
	
	The formal possibility of reconstructing $C_\pi^{(r)}(t)$ in the standard
	formulation---by reading off the explicit $8\times8$ matrix for $P_r$ from the
	$\cl(6)$ calculation and using it as an ad hoc operator---does not constitute
	an independent standard QCD construction.  The matrix would be algebraically
	unmotivated without $H$; its gauge invariance $[T_a, P_r]=0$ would require
	separate verification; and its physical interpretation as a projector onto a
	specific $\su(3)$ representation subspace of the spinor would be inaccessible.
	Moreover, the $\cl(6)$ quark space
	$\mathbb{C}^{24} = \mathbb{C}^8\otimes\mathbb{C}^3$ does not embed naturally
	into the standard quark space $\mathbb{C}^{12} = \mathbb{C}^4\otimes\mathbb{C}^3$:
	the factor-of-two mismatch in spinor dimension means $P_r$ cannot be directly
	expressed as a matrix in the $12\times12$ space without first specifying an
	embedding of $\mathbb{C}^4$ into $\mathbb{C}^8$, a choice that is itself
	non-canonical and requires the $\cl(6)$ structure to motivate.
	
	The $\cl(6)$ formulation provides three things that the standard formulation
	cannot: the internal generators $e_5, e_6$ that make $H$ and $P_r$ definable;
	the algebraic guarantees---$[T_a, P_r]=0$ (gauge invariance),
	$[\gamma_5, P_r]=0$ (chiral compatibility), $\sum_r P_r = I_8$
	(completeness)---all following automatically from the algebra without separate
	proof; and the $(\GaSeven, H)$ interpretation that identifies the $\su(3)$
	representation content of each channel.  The grade decomposition of the pion
	correlator is accordingly not a renaming of a standard QCD observable, but a
	structure that is natural and algebraically guaranteed in $\cl(6)$ and has no
	natural standard QCD expression.
	
	It is worth distinguishing two senses in which this structure might be called
	novel.  The grade decomposition is \emph{definable} only in $\cl(6)$: the
	construction above requires the internal generators $e_5,e_6$ and is independent
	of any dynamics, so in that structural sense it is not a relabeling of a standard
	color projection.  Whether it additionally \emph{carries} physical information
	inaccessible to standard QCD is a distinct question, and one the present
	heavy-mass data do not settle: at $m_0=0.05$ the measured fractions remain
	pinned to their kinematic values $\dim(r)/8$ (Sec.~\ref{sec:experiment-A}), so
	at this parameter point the decomposition exhibits no resolvable dynamical
	content.  Testing the informational sense---whether the grade fractions resolve
	gauge dynamics not already captured by ordinary color and chiral
	observables---is the purpose of Experiment~B
	(Sec.~\ref{sec:experiment-B-design}).
	
	Proposition~\ref{prop:gamma7-frozen} sharpens this question and bounds it
	from one side.  The $\spin(6)$-chirality direction---the $(e_5,e_6)$ content
	carried by $\GaSeven$---is \emph{provably} information-free:
	$f_{(\GaSeven=+1)} = f_{(\GaSeven=-1)} = \half$ on every configuration,
	independent of $\beta$, of $t$, and of the gauge field, precisely because
	$e_5,e_6$ never enter $\dW$.  Whatever dynamical content the grade fractions
	can ever resolve must therefore reside entirely in the splitting \emph{within}
	a fixed $\GaSeven$ sector---the triplet-versus-singlet $H$-splitting of the
	$(\GaSeven,H)$ double grading, a $\gamma_5$-type four-dimensional structure
	dressed by color.  This is exactly the single C-even observable
	$\ftriplet = f_3 + f_{\bar 3}$ isolated in Sec.~\ref{sec:channel-pair}, and it
	is the target of Experiment~B.  The exact identity thus does not weaken the
	novelty claim: it removes a direction in which no new information could ever
	appear, and concentrates the informational question onto the one observable
	capable of carrying it.
	
	
	
	\subsection{Outlook}
	\label{sec:outlook}
	
	\subsubsection{The random-walk hypothesis for the chiral regime}
	\label{sec:random-walk-hypothesis}
	
	The algebraic result~\eqref{eq:H-e1-commutator}, $[H,e_\mu]\neq0$, implies
	that every spacetime hop of the quark shifts its $H$-eigenvalue.  More
	precisely, each application of $e_\mu$ for $\mu\in\{1,2,3,4\}$ acts as a
	$\pm1$ step on the $H$-eigenvalue ladder
	$\{-\tfrac{3}{2},\,-\tfrac{1}{2},\,+\tfrac{1}{2},\,+\tfrac{3}{2}\}$, and
	the accumulated product $G=\dW^{-1}$ therefore constitutes a random walk in
	the grade space of the eight-dimensional spinor.  The stationary distribution
	of this walk---the long-time limit of the grade fraction $f_r(t)$---is
	determined by the gauge dynamics, not by free-field counting.
	
	At heavy quark mass ($m_0=0.05$, $m_{\mathrm{phys}}\approx0.76$), the
	propagator is short-range and the walk effectively terminates after $O(1)$
	hops.  The gauge-link weighting of each step is small ($|U_\mu-I_3|$ is
	perturbative in $g^2$), the mixing is nearly symmetric, and the free-field
	theorem controls the grade fractions: $f_r\approx\dim(r)/8$ at the $10^{-4}$
	level, as confirmed by Experiment~A.
	
	At $\kappa\to\kcren$, the pion mass $m_\pi\to0$ and the propagator becomes
	long-range.  For $t\,m_\pi\ll1$ the kinematic regime persists and the
	free-field theorem still governs the grade fractions.  For $t\,m_\pi\sim1$,
	however, the propagator has traversed $O(m_\pi^{-1})$ lattice spacings and
	accumulated $O(m_\pi^{-1})$ gauge-link-weighted $H$-mixing steps.  The grade
	fractions $f_r(t)$ will then curve from their free-field values toward an
	asymptotic limit determined by the ground-state pion wavefunction in the grade
	space: $f_r(t\to\infty)\to Z_\pi^{(r)}/Z_\pi$.
	
	The central hypothesis motivating Experiment~B is that the confined gauge
	vacuum ($\beta=4.5$) biases this walk toward the color-triplet $H$-eigenspaces
	more strongly than the deconfined vacuum ($\beta=6.0$), producing a nonzero
	phase difference
	\begin{equation}
		\dftriplet \;\equiv\;
		\ftriplet(\beta=4.5) \;-\; \ftriplet(\beta=6.0) \;\neq\; 0
		\label{eq:delta-f-triplet}
	\end{equation}
	near $\kcren$.  The motivating argument has two parts, only the first of
	which is established.
	
	The first part is a standard lattice-QCD expectation, not a measurement made
	here.  In the chiral regime that Experiment~B would probe, the strongly
	disordered gauge vacuum of the confined phase is expected to support a denser
	spectrum of near-zero Dirac modes than the deconfined phase: by analogy with
	the Banks--Casher mechanism $\Sigma=\pi\rho(0)$~\cite{BanksCasher1980}, in
	which the chiral condensate $\Sigma$ is proportional to the spectral density
	at the origin, $|\Sigma|$ would be larger at $\beta=4.5$ than at $\beta=6.0$.
	Such low modes would dominate the long-range pion propagator near $\kcren$
	through the spectral decomposition
	$G(x,t;0,0)=\sum_n\psi_n(x,t)\psi_n^\dagger(0,0)/\lambda_n$, in which the
	small-$\lambda_n$ contributions are weighted most heavily.  This chain is
	conventional lattice QCD; the present heavy-mass work neither reaches that
	regime nor measures $\rho(0)$.
	
	The second part is conjectural and is precisely what Experiment~B is designed
	to test: that the low Dirac modes, when projected onto the $\cl(6)$ spinor
	space, carry preferential weight in the color-triplet $H$-eigenspaces
	($H=\pm1/2$) relative to the singlet $H$-eigenspaces ($H=\pm3/2$).  This
	$H$-grade preference, if present, would propagate through the spectral sum to
	produce the bias~\eqref{eq:delta-f-triplet}.  In standard QCD the
	Atiyah--Singer index theorem fixes the $\gamma_5$ chirality of zero modes in
	nontrivial topological sectors but does not constrain any analog of the
	$H$-grade content.  No standard QCD argument predicts a preference between the
	$H=\pm1/2$ and $H=\pm3/2$ eigenspaces of the $\cl(6)$ Cartan hypercharge, and
	we are not aware of any prior work that addresses this question.  We offer the
	hypothesis as a physically motivated guess, with the explicit acknowledgment
	that the opposite sign (singlet enhanced) and the null result ($\dftriplet=0$)
	are equally consistent with present knowledge.
	
	We stress that~\eqref{eq:delta-f-triplet} is a hypothesis, not a derived
	result.  Both the magnitude and the sign of $\dftriplet$ require the
	Experiment~B data to confirm or refute.  The free-field theorem provides the
	null hypothesis: $\dftriplet=0$.  Experiment~A establishes this null at heavy
	mass ($0.6\sigma$); Experiment~B will test it at the chiral point.
	
	\subsubsection{Experiment~B: design and falsification criteria}
	\label{sec:experiment-B-design}
	
	Experiment~B will run at $\kappa\approx\kcren(\beta)$ determined from Block~1:
	$\kappa\approx0.143$--$0.145$ at $\beta=4.5$ and $\kappa\approx0.150$ at
	$\beta=6.0$ (systematic uncertainty $\pm0.005$ from the $N=4$/$N=6$ spread;
	see Sec.~\ref{sec:pcac-method}).  The target observable is the C-even phase
	difference $\dftriplet$ defined in Eq.~\eqref{eq:delta-f-triplet}.  The three
	practical complications
	identified in Sec.~\ref{sec:checkerboard} require specific mitigations before
	running: ($i$) CG ill-conditioning near $\kcren$ is addressed by deflation of
	the ten lowest eigenmodes, even--odd Schur complement preconditioning, and
	adaptive solver tolerance set to $10^{-10}$, as detailed in
	Appendix~\ref{app:D}; ($ii$) signal-to-noise degradation at large $t$ restricts
	the analysis to timeslices satisfying $C_\pi(t)/\sigma_{C_\pi}(t)>3$ in both
	phases; ($iii$) the $\kcren$ systematic is controlled by running at two values
	$\kcren\pm0.003$ and verifying that $\dftriplet$ is stable across this range.
	
	The falsification criteria are as follows.  A positive outcome---$\dftriplet>0$
	at $\geq3\sigma$ significance with $\Ncfg\gtrsim100$ at $N=6$---would
	constitute evidence for the hypothesis: that the chiral gauge vacuum
	distributes spinor weight preferentially into the color-carrying channels in
	the confined phase.  Confirmation, in the lattice-QCD sense of an effect that
	survives systematic checks, would require independent verification by volume
	scaling to $N=8$ and a continuum-limit study at multiple $\beta$
	(Sec.~\ref{sec:diagnostic-roadmap}); the present work cannot make a
	confirmation claim on a single volume at a single lattice spacing regardless of
	per-point significance.  A null outcome---$\dftriplet$ consistent with zero at
	$\Ncfg\gtrsim100$---would establish that grade fractions are insensitive to the
	confinement--deconfinement transition even at the chiral point, and that the
	free-field theorem extends to the long-range propagator regime; this result
	would itself be clean and publishable, establishing the grade fractions as a
	renormalization-free null observable.  An anti-screening
	outcome---$\dftriplet<0$---would be physically unexpected and would prompt
	investigation of whether the effect is genuine (requiring volume scaling) or
	numerical (CG convergence bias near $\kcren$).  Persistence of
	$A_{\mathrm{color}}$ or $A_{\mathrm{singlet}}$ beyond the $\Ncfg^{-1/2}$
	scaling at large statistics would indicate a C-symmetry violation in the code
	or ensemble update and must be diagnosed before any physics interpretation is
	attempted.
	
	\subsubsection{Diagnostic roadmap}
	\label{sec:diagnostic-roadmap}
	
	Before any physical interpretation of Experiment~B results, a free-field
	prerequisite test must be performed.  With all gauge links set to
	$U_\mu(x)=I_3$, the projector code, the trace evaluation, and the propagator
	pipeline must reproduce $f_r=\dim(r)/8$ to machine precision on every
	configuration, in accordance with Theorem~\ref{thm:kinematic}.  This requires
	no gauge generation and runs in minutes on a single workstation.  It is a
	prerequisite for trusting any nonzero result from the chiral-regime
	calculation: the test was verified for the heavy-mass production code
	(Sec.~\ref{sec:experiment-A}, where $\delta_{\max}=5.4\times10^{-16}$ was
	observed) and must be re-verified after any code modifications introduced for
	the solver mitigations of Sec.~\ref{sec:experiment-B-design}.  Failure
	indicates a bug in the projector or trace pipeline and invalidates any
	subsequent physics interpretation; passing is necessary but not sufficient for
	the chiral-regime results to be physical.
	
	Beyond this prerequisite, the three mechanisms identified in
	Sec.~\ref{sec:deviation-interp} can be disentangled by a sequence of targeted
	calculations ordered by computational cost.
	
	The first physics test is $\Ncfg$ scaling: increasing the ensemble size from
	$\Ncfg=20$ to $\Ncfg\approx200$ at the current heavy-mass parameters
	($\kappa=0.12346$, $\beta\in\{4.5,6.0\}$, $N=6$).  If per-channel deviations
	from $\dim(r)/8$ scale as $\Ncfg^{-1/2}$ and approach zero, mechanism (a)
	is confirmed as the primary driver and no floor from (b) or (c) is present.
	If deviations approach a nonzero floor, the sign and structure would
	distinguish a lattice artifact (mechanism b, positive-definite) from a
	vacuum-dressing signal (mechanism c, phase-dependent).
	
	Separating mechanisms (b) and (c) rigorously requires a continuum-limit study:
	repeating Experiment~B at multiple values of $\beta$ while holding the physical
	pion mass fixed by adjusting $\kappa$ accordingly.  Mechanism (b) contributes
	at $O(a)$ and vanishes as $\beta\to\infty$; mechanism (c) is physical and
	survives the continuum limit.  A practical program would use
	$\beta\in\{4.5,\,5.5,\,6.0,\,6.5\}$ with corresponding $\kcren$ values from
	Block~1 (Sec.~\ref{sec:pcac-method}), on $N=6$ or $N=8$ lattices.  Concretely,
	this requires three extensions beyond the present work: Block~1 must be
	completed to all six $\beta$ values; the Wilson action should be improved to
	remove leading $O(a)$ artifacts; and ensembles must be generated at multiple
	lattice spacings.  Collectively this is a project at least an order of
	magnitude larger than the present work and is beyond the single-workstation
	budget within which this program currently operates.
	
	\subsubsection{Future directions}
	\label{sec:future-directions}
	
	Two directions extend the grade-decomposition program beyond what
	Sec.~\ref{sec:diagnostic-roadmap} addresses.  Both lie outside the scope of
	the present work but are noted as natural targets for follow-up investigation.
	
	The first is the extension to dynamical fermions.  In the quenched
	approximation used throughout, the fermion-loop contributions carrying the
	$\mathrm{U}(1)_A$ anomaly are suppressed by construction; the singlet
	subspaces $\mathbf{1}_a$ and $\mathbf{1}_b$ are therefore insensitive to
	topological charge at the present level of approximation.  Unquenching could couple the axial anomaly to the grade channels through
	near-zero Dirac modes, but Proposition~\ref{prop:gamma7-frozen} sharply
	limits where such an effect can appear.  The $\spin(6)$-chirality combination
	$A_{\GaSeven} = A_{\mathrm{color}} + A_{\mathrm{singlet}}$ vanishes identically
	on every configuration because $e_5,e_6$ never enter $\dW$; unquenching
	modifies the gauge \emph{measure} but not this operator content, so the
	identity---and with it the per-configuration lock
	$A_{\mathrm{singlet}} = -A_{\mathrm{color}}$ (equivalently $f_{1_a}=\half-f_3$,
	$f_{1_b}=\half-f_{\bar 3}$)---survives unquenching unchanged.  The singlet
	asymmetry therefore cannot acquire a topological-sector dependence
	\emph{independent} of the color asymmetry: any topological response of
	$A_{\mathrm{singlet}}$ is exactly minus that of $A_{\mathrm{color}}$, and the
	chirality direction as a whole carries none.  The only grade observable in
	which dynamical-fermion or anomaly effects could register is the
	within-$\GaSeven$ $H$-splitting $\ftriplet = f_3 + f_{\bar 3}$
	(Sec.~\ref{sec:cl6-specific})---equivalently the complementary singlet total
	$f_{1_a}+f_{1_b}=1-\ftriplet$.  Whether the $\eta'$ anomaly leaves a measurable
	imprint there is a speculative connection requiring dynamical-fermion ensembles
	and explicit topological-charge measurement to test.
	
	The second is the connection to geometric-algebra-equivariant machine learning.
	The four grade-projected correlators $\{C_\pi^{(r)}(t)\}$ share properties
	that make them a natural target for GA-equivariant neural network architectures:
	they are gauge-invariant by construction ($[T_a,P_r]=0$), mutually orthogonal
	($P_rP_s=0$ for $r\neq s$), and labeled by a discrete quantum number---the
	$(\GaSeven,H)$ pair---encoding the $\su(3)$ representation content of each
	channel.  Clifford-equivariant
	networks~\cite{Ruhe2023a,Brehmer2023,Brehmer2024,Ruhe2024} are designed to
	respect exactly this type of graded algebraic structure as an inductive bias,
	and the double grading of Sec.~\ref{sec:grading-literature} provides a
	ready-made label set---triplet, antitriplet, singlet-$a$, singlet-$b$---for
	supervised or unsupervised decomposition of lattice gauge data.  Whether this
	algebraic compatibility translates into practical advantage in generalization
	or data efficiency remains an open question that the present work does not
	pursue; we note the connection as a possible application of the $\cl(6)$
	framework to geometric-algebra-equivariant deep learning for physical systems,
	and defer any concrete claims to future investigation.
		
		Finally, the broader programme deferred in Sec.~\ref{sec:scope} remains open as
		independent follow-up: a renormalized Gell-Mann--Oakes--Renner analysis with
		Symanzik-improved fermions to control the additive Wilson renormalization that
		obscures it here, and a direct measurement of the infrared spectral density
		$\rho(0)$ in the chiral regime to replace the heuristic Banks--Casher argument
		of Sec.~\ref{sec:random-walk-hypothesis} with a computed quantity.  A second,
		independent determination of $\kcren$---for instance via the $m_\pi^2$
		extrapolation that proved unstable at $N\le6$ (Sec.~\ref{sec:pcac-method})---would
		also supply the cross-check that the present single-method value lacks.
	
	% ============================================================
	\section{Conclusion}
	\label{sec:conclusion}
	% ============================================================
	
	We have constructed and numerically implemented the grade decomposition of
	the pion correlator in the Clifford algebra $\cl(6)$ framework.  The
	principal findings are as follows.
	
	\noindent\textbf{Algebraic construction.}  The four-channel decomposition
	$C_\pi(t) = \sum_r C_\pi^{(r)}(t)$ follows from the Lagrange spectral
	projectors of the Cartan hypercharge $H = B_{12}+B_{34}+B_{56}$ of
	$\mathrm{spin}(6)$.  The projectors satisfy six independent algebraic
	identities---completeness, idempotency, orthogonality, $\su(3)$ covariance,
	charge-conjugation pairing, and $\gamma_5$-commutativity---all verified to
	machine precision $\leq 2.4\times10^{-15}$.  The $(\GaSeven, H)$ double
	grading is the standard $\su(4)\supset\su(3)\times\mathrm{U}(1)$ Levi
	decomposition; the present paper is its first dynamical implementation.
	
	\noindent\textbf{Kinematic limit theorem.}  In the free Wilson--Dirac theory,
	the grade fractions $f_r(t) = \dim(r)/8$ exactly on every gauge
	configuration.  The proof proceeds through the $I_4\otimes\mathrm{U}(2)$
	symmetry of the free propagator, Schur's lemma, free-field color triviality,
	and the spatial-cubic invariance of the spatially-summed correlator---the
	residual time-direction term being removed by its vanishing grade-basis
	diagonal (Theorem~\ref{thm:kinematic}).  The theorem is verified
	at $\delta_{\max} = 1.26\times10^{-16}$---machine precision---on the free
	lattice.
	
	\noindent\textbf{Experiment A: heavy-mass production.}  On thermalized
	quenched ensembles at $N=6$, $m_0=0.05$, $\beta\in\{4.5,6.0\}$,
	$\Ncfg=20$: the algebraic completeness identity holds to
	$\delta_{\max}=5.4\times10^{-16}$ on every configuration; the measured
	grade fractions are consistent with the kinematic free-field values at the
	$10^{-4}$ level; and the phase differences are consistent with zero at
	$0.6\sigma$.  These results establish that the grade decomposition machinery
	works at machine precision and that the heavy-mass result is consistent with
	the kinematic theorem.  They do not address the chiral regime.
	
	\noindent\textbf{Scope statement.}  Three limitations apply: single lattice
	volume ($N=6$), single lattice spacing, and quenched approximation.  The
	central physical question---whether non-perturbative gauge dynamics at the
	chiral point distinguish the representation channels---requires Experiment~B
	(Sec.~\ref{sec:experiment-B-design}), which is identified as the essential
	follow-up.
	
	\noindent\textbf{Structural observation.}  The grade decomposition has no
	natural analog in the standard four-component Wilson--Dirac formulation: the
	bivector $B_{56}$~\eqref{eq:B56} requires the internal generators $e_5, e_6$
	of $\cl(6)$ that are absent from $\cl(4)$, and the grade structure cuts
	across the $\gamma_5$ chirality decomposition.  The observables
	$C_\pi^{(r)}(t)$ are accordingly $\cl(6)$-specific; their gauge invariance,
	completeness, and $\su(3)$ representation-theoretic interpretation all follow
	automatically from the algebraic structure without separate proof.  A sharper
	structural fact further constrains what they can measure: because $e_5,e_6$ do
	not enter $\dW$, the $\spin(6)$-chirality fractions are frozen at
	$f_{(\GaSeven=\pm1)}=\half$ on every configuration
	(Proposition~\ref{prop:gamma7-frozen}), so the only dynamically resolvable
	content resides in the within-$\GaSeven$ $H$-splitting
	$\ftriplet = f_3 + f_{\bar 3}$---the observable Experiment~B is designed to
	probe.
	
	Paper~1~\cite{FirestoneP1} established the $\cl(6)$ formulation as a valid
	and computationally tractable framework for Wilson lattice QCD.  The present
	paper establishes the algebraic infrastructure for grade-decomposed observables
	and confirms the kinematic null result at heavy mass.  Experiment~B will
	determine whether the grade decomposition is a $\cl(6)$-specific renaming of
	the standard result or carries genuinely new physical information.
	
	\medskip
	\noindent\textbf{Data and code availability.}  All simulation code, raw data,
	and analysis scripts are publicly available under the MIT license.
	% NOTE: Insert repository URL (github.com/...) before submission.
	Environment: Python~3.11.9, NumPy~2.2.5,
	SciPy~$\geq$~1.11.
	
	% ============================================================
	\appendix
	\renewcommand{\thetable}{\Alph{section}.\arabic{table}}
	% ============================================================
	
	% ============================================================
	\section{\texorpdfstring{$\cl(6)$}{Cl(6)} Basis and Grade Structure}
	\label{app:A}
	% ============================================================
	
	The $\cl(6)$ basis and multiplication structure are presented in full in
	Paper~1~\cite{FirestoneP1} Appendix~A.  Here we collect only the elements
	most directly relevant for the grade projector construction of
	Sec.~\ref{sec:projectors}.
	
	The 15 bivectors $B_{ij}=(i/4)[e_i,e_j]$ are all Hermitian and traceless,
	with eigenvalues exactly $\{-\tfrac{1}{2}\ (\times4),\,+\tfrac{1}{2}\ (\times4)\}$.
	The Cartan hypercharge of the $\su(4)\supset\su(3)\times\mathrm{U}(1)$ Levi
	decomposition is $H = B_{12}+B_{34}+B_{56}$.  In the tensor product
	representation~\eqref{eq:basis},
	\[
	H \;=\; -\tfrac{1}{2}(\sigma_z\otimes I_2\otimes I_2)
	- \tfrac{1}{2}(I_2\otimes\sigma_z\otimes I_2)
	- \tfrac{1}{2}(I_2\otimes I_2\otimes\sigma_z),
	\]
	which is diagonal with entries
	$\bigl(-\tfrac{3}{2},\,-\tfrac{1}{2},\,-\tfrac{1}{2},\,+\tfrac{1}{2},\,
	-\tfrac{1}{2},\,+\tfrac{1}{2},\,+\tfrac{1}{2},\,+\tfrac{3}{2}\bigr)$
	in tensor product order, consistent with~\eqref{eq:H-diag} and the
	$(\GaSeven,H)$ double grading of Sec.~\ref{sec:double-grading}.
	
	The pseudoscalar $\Isix = e_1e_2e_3e_4e_5e_6$ satisfies $\Isix^2 = -1$, is anti-Hermitian, and
	anticommutes with all odd-grade elements of $\cl(6)$.  The $\spin(6)$
	chirality $\GaSeven = i\Isix$ is Hermitian with eigenvalues $\pm1$ and
	commutes with $H$ and with every projector $P_r$.  The 4D chiral operator
	$\gamma_5 = e_1e_2e_3e_4$ satisfies $\gamma_5^2 = I_8$, anticommutes with
	$e_1,\ldots,e_4$, and commutes with $e_5$, $e_6$, $H$, and all $P_r$
	(verified to $0.00\times10^0$ in Sec.~\ref{sec:proj-gamma5}).  The
	distinction $\GaSeven \neq \gamma_5$ is central to the physics interpretation
	of Sec.~\ref{sec:cl6-specific}.
	
	% ============================================================
	\section{Numerical Verification Details}
	\label{app:B}
	% ============================================================
	
	\noindent\textbf{Environment:} Python~3.11.9 $\cdot$ NumPy~2.2.5 $\cdot$
	SciPy~$\geq$1.11 $\cdot$ IEEE~754 double precision
	(machine $\varepsilon \approx 2.22\times10^{-16}$)
	% NOTE: Insert repository URL (github.com/...) before submission.
	
	The foundational $\cl(6)$ and $\su(3)$ algebraic verifications are reprised
	from Paper~1~\cite{FirestoneP1} Appendix~B without change.  Table~B1 of that
	appendix records the $\cl(6)$ gamma matrix properties; Table~B2 the $\su(3)$
	generator properties (maximum commutator error $1.24\times10^{-16}$); Table~B5
	the coupling structure.  Full details are given in
	Paper~1~\cite{FirestoneP1} Appendix~B.
	
	New algebraic verifications introduced in this paper are collected in the
	table below.  All residuals are maximum Frobenius norms over the relevant
	index set.  The final block verifies that the bivector $\su(3)$ octet,
	derived as the stabilizer of the preferred spinor $|\xi_0\rangle$ (the
	$H=-\tfrac{3}{2}$ eigenstate) within the 15-dimensional bivector space,
	closes as a genuine nonabelian Lie subalgebra of $\spin(6)$, and that
	its centralizer in $\mathrm{span}(B_{ij})$ is exactly
	$\mathrm{span}\{H\}$.
	
	\medskip
	\begin{center}
		\begin{tabular}{@{}lcc@{}}
			\toprule
			\textbf{Check} & \textbf{Target} & \textbf{Result} \\
			\midrule
			$H$ eigenvalue spectrum:
			$\bigl\{-\tfrac{3}{2},-\tfrac{1}{2},+\tfrac{1}{2},+\tfrac{3}{2}\bigr\}$,
			multiplicities $\{1,3,3,1\}$
			& Exact & $\checkmark$\ Exact \\[6pt]
			$[\GaSeven,\,H] = 0$
			& $0.00\times10^{0}$ & $\checkmark$\ $0.00\times10^{0}$ \\[6pt]
			$[H,\,T_a] = 0$,\; all $a\in\{1,\ldots,8\}$
			& $<10^{-15}$ & $\checkmark$\ $<10^{-15}$ \\[6pt]
			$\gamma_5$ commutativity:\;
			$[\gamma_5,\,P_r] = 0$,\; all $r$
			& $0.00\times10^{0}$ & $\checkmark$\ $0.00\times10^{0}$ \\[6pt]
			Charge conjugation:\;
			$C P_{\mathbf{3}} C^{\dagger} = P_{\bar{\mathbf{3}}}$,\;
			$C P_{\mathbf{1}_a} C^{\dagger} = P_{\mathbf{1}_b}$
			& $0.00\times10^{0}$ & $\checkmark$\ $0.00\times10^{0}$ \\[6pt]
			$\Isix P_{\mathbf{1}_a} = -i\,P_{\mathbf{1}_a}$,\;
			$\Isix P_{\mathbf{1}_b} = +i\,P_{\mathbf{1}_b}$
			& $0.00\times10^{0}$ & $\checkmark$\ $0.00\times10^{0}$ \\[6pt]
			\multicolumn{3}{@{}l}{\textit{Bivector $\su(3)$ octet
			(stabilizer of $|\xi_0\rangle$, $H=-\tfrac{3}{2}$):}} \\[3pt]
			\quad Stabilizer dimension in $\mathrm{span}(B_{ij})$
			& $8$ & $\checkmark$\ $8$ \\[6pt]
			\quad Lie-algebra closure:\;
			$\max\bigl\|[X_a,X_b]-\mathrm{proj}_{\su(3)}[X_a,X_b]\bigr\|_F$
			& $<10^{-14}$ & $\checkmark$\ $6.2\times10^{-16}$ \\[6pt]
			\quad Nonabelian:\;
			$\max\|[X_a,X_b]\|_F$
			& $>0$ & $\checkmark$\ $2.0$ \\[6pt]
			\quad Centralizer of $\su(3)$ in $\mathrm{span}(B_{ij})$:\;
			$\dim=1$,\; $\bigl|\langle\hat{v},\,\hat{H}\rangle\bigr|$
			& $1.0$ & $\checkmark$\ $1.0000000000$ \\
			\bottomrule
		\end{tabular}
	\end{center}
	\medskip
	
	\noindent\textbf{Checkerboard sweep correctness.}\enspace
	The corrected even-odd Metropolis sweep produces equilibrium plaquettes
	within $0.7\%$ of the Paper~1~\cite{FirestoneP1} reference values at $N=6$
	(Table~\ref{tab:plaquette-gate}), consistent with the residual difference
	being attributable to the parallel-sweep systematic in earlier intermediate
	runs; see Sec.~\ref{sec:checkerboard} for the methodological correction.
	
	% ============================================================
	\section{Grade Projector Verification}
	\label{app:C}
	% ============================================================
	
	Table~\ref{tab:verif-C1} reports the six projector verification checks
	referenced in Secs.~\ref{sec:proj-construction}
	and~\ref{sec:proj-verification}.  All checks were performed in IEEE~754
	double precision, Python~3.11.9, NumPy~2.2.5.  Checks C1--C5 must pass
	to machine precision before any production computation proceeds; they are
	algebraic prerequisites, not optional verifications.
	
	\setcounter{table}{0}
	\begin{table}[htbp]
		\centering
		\caption{Grade projector algebraic verification.}
		\label{tab:verif-C1}
		\begin{tabular}{@{} c p{6.3cm} c c c @{}}
			\toprule
			\textbf{Check}
			& \textbf{Identity}
			& \textbf{Target}
			& \textbf{Result}
			& \textbf{Status} \\
			\midrule
			C1
			& Completeness:\; $\max\bigl\|\sum_r P_r - I_8\bigr\|_F$
			& $<10^{-14}$
			& $0.00\times10^{0}$
			& $\checkmark$\ PASS \\[4pt]
			C2
			& Idempotency:\; $\max\|P_r^2 - P_r\|_F$
			& $<10^{-14}$
			& $0.00\times10^{0}$
			& $\checkmark$\ PASS \\[4pt]
			C3
			& Orthogonality:\; $\max\|P_r P_s\|_F$,\; $r\neq s$
			& $<10^{-14}$
			& $0.00\times10^{0}$
			& $\checkmark$\ PASS \\[4pt]
			C4
			& Dimensions: $\tr(P_{\mathbf{3}})=3$,\; $\tr(P_{\bar{\mathbf{3}}})=3$,\;
			$\tr(P_{\mathbf{1}_a})=1$,\; $\tr(P_{\mathbf{1}_b})=1$
			& exact integers
			& exact integers
			& $\checkmark$\ PASS \\[4pt]
			C5
			& $\su(3)$ covariance: $\max\|U P_r U^\dagger - P_r\|_F$
			over 50 random $U\in\su(3)$, all $r$
			& $<10^{-12}$
			& $2.44\times10^{-15}$
			& $\checkmark$\ PASS \\[4pt]
			S1$^\dagger$
			& Haar average:
			$\langle\rho(g)\rangle \to P_{\mathbf{1}_a}+P_{\mathbf{1}_b}$;\;
			$N_{\mathrm{samples}}=50000$ random $U\in\su(3)$
			& ${\sim}N_{\mathrm{samples}}^{-1/2}$
			& $0.0050$\newline($1/\!\sqrt{N}$ OK)
			& $\checkmark$\ PASS \\
			\bottomrule
		\end{tabular}
		
		\smallskip
		\noindent\footnotesize
		$^\dagger$\ Stochastic check.  The algebraic identity is established by
		C1--C5; S1 verifies behavior under finite-sample group averaging.  Expected
		error $\sigma_{\max}/\sqrt{N}\approx0.0026$; observed $0.0050$; the ratio
		$1.9$ is of the order of, and bounded by, the max-over-element estimate
		$\sqrt{2\ln 64}\approx 2.9$.
		The $1/\sqrt{N}$ scaling is verified across $N_{\mathrm{samples}}=100$ to
		$200{,}000$.
	\end{table}
	
	% ============================================================
	\section{CG Solver Mitigations for the Chiral Regime}
	\label{app:D}
	% ============================================================
	
	The heavy-mass production runs ($m_0=0.05$, $\kappa=0.12346$) required no
	special solver treatment: mean CG iteration counts were $18.95$ at
	$\beta=4.5$ and $21.85$ at $\beta=6.0$, with zero solver failures across
	all $\Ncfg=20$ configurations.  This appendix provides the prerequisite
	recipe for Experiment~B (Sec.~\ref{sec:experiment-B-design}) at
	$\kappa\approx\kcren$, where the CG condition number
	$\kappa_{\mathrm{cond}}\sim 1/m_{\mathrm{phys}}^2$ diverges.  The three
	mitigations below are recommended in order of implementation priority.
	
	\medskip
	\noindent\textbf{(i) Deflation of low modes.}\enspace
	Compute the 10 lowest eigenmodes of $\dW^\dagger\dW$ using LOBPCG or a
	Lanczos routine before the main CG sweeps.  Subtract the deflated
	contribution exactly before applying CG; this removes the conditioning
	bottleneck from the near-zero sector.  Expected speedup: $5$--$10\times$
	near $\kcren$ on $N=6$ lattices.
	
	\medskip
	\noindent\textbf{(ii) Even-odd Schur complement preconditioning.}\enspace
	Decompose the lattice into even and odd sites.  The Schur complement
	$D_{ee} - D_{eo}D_{oo}^{-1}D_{oe}$ has a better condition number than
	$\dW$ itself at strong coupling.  $D_{oo}$ is diagonal in the hopping
	parameter $\kappa$ representation and invertible exactly.  This reduces the
	effective system size by a factor of two and removes the odd-site
	contribution to conditioning.
	
	\medskip
	\noindent\textbf{(iii) Adaptive tolerance.}\enspace
	Within $5\%$ of $\kcren$ (estimated from Table~\ref{tab:kappa-c}), tighten
	the CG stopping criterion to $\mathrm{rtol}=10^{-10}$ and raise
	\texttt{max\_iter} to 5000.  Monitor the residual
	$\|\dW x - b\|/\|b\|$ post-solve; if it exceeds $10^{-3}$, discard the
	configuration and log the failure rather than proceeding with a spurious
	propagator.
	
	\medskip
	If CG still fails to converge on $N=10$ lattices after applying all three
	mitigations, this is a hard limitation to be reported honestly.  No results
	should be produced by forcing convergence.  Experiment~B is not required to
	succeed on $N=10$ for the null-or-signal outcome to be interpretable on
	$N=6$.
	
	% ============================================================
	%  BIBLIOGRAPHY
	%  Manual thebibliography — do NOT run BibTeX or replace with .bib.
	%  Labels [20]–[28] enforced via \bibitem[N]{key} to preserve
	%  continuity with Paper 1 reference numbering.
	%  [P1] uses custom label \bibitem[{[P1]}]{FirestoneP1}.
	%  NOTE TO EDITOR: reference numbers [12]--[19] are intentionally reserved so
	%  that citations shared with Paper 1 retain identical numbers; the resulting
	%  non-contiguous numbered list is deliberate, not an omission.  (This note is
	%  also included in the submission cover letter.)
	% ============================================================
	
	\begin{thebibliography}{99}
		
		%% [P1] — Paper 1 (custom label; not part of numbered sequence)
		\bibitem[{[P1]}]{FirestoneP1}
		T.~M. Firestone,
		``Clifford algebra $\cl(6)$ formulation of lattice QCD: algebraic structure,
		Dirac operator construction, and verification against standard benchmarks,''
		% NOTE: Insert publication details (journal / arXiv ID / year) when available.
		
		%% [1]–[11]: standard numbered references
		\bibitem[1]{PatiSalam1974}
		J.~C. Pati and A.~Salam,
		``Lepton number as the fourth color,''
		\textit{Phys.\ Rev.}\ D~\textbf{10}, 275 (1974).
		\href{https://doi.org/10.1103/PhysRevD.10.275}{doi:10.1103/PhysRevD.10.275}.
		
		\bibitem[2]{Slansky1981}
		R.~Slansky,
		``Group theory for unified model building,''
		\textit{Phys.\ Rep.}\ \textbf{79}, 1 (1981).
		\href{https://doi.org/10.1016/0370-1573(81)90092-2}{doi:10.1016/0370-1573(81)90092-2}.
		
		\bibitem[3]{BanksCasher1980}
		T.~Banks and A.~Casher,
		``Chiral symmetry breaking in confining theories,''
		\textit{Nucl.\ Phys.}\ B~\textbf{169}, 103 (1980).
		\href{https://doi.org/10.1016/0550-3213(80)90255-2}{doi:10.1016/0550-3213(80)90255-2}.
		
		\bibitem[4]{Wilson1974}
		K.~G. Wilson,
		``Confinement of quarks,''
		\textit{Phys.\ Rev.}\ D~\textbf{10}, 2445 (1974).
		\href{https://doi.org/10.1103/PhysRevD.10.2445}{doi:10.1103/PhysRevD.10.2445}.
		
		\bibitem[5]{KogutSusskind1975}
		J.~B. Kogut and L.~Susskind,
		``Hamiltonian formulation of Wilson's lattice gauge theories,''
		\textit{Phys.\ Rev.}\ D~\textbf{11}, 395 (1975).
		\href{https://doi.org/10.1103/PhysRevD.11.395}{doi:10.1103/PhysRevD.11.395}.
		
		\bibitem[6]{Gattringer2010}
		C.~Gattringer and C.~B. Lang,
		\textit{Quantum Chromodynamics on the Lattice}
		(Springer, Berlin, 2010).
		\href{https://doi.org/10.1007/978-3-642-01850-3}{doi:10.1007/978-3-642-01850-3}.
		
		\bibitem[7]{Sheikholeslami1985}
		B.~Sheikholeslami and R.~Wohlert,
		``Improved continuum limit lattice action for QCD with Wilson fermions,''
		\textit{Nucl.\ Phys.}\ B~\textbf{259}, 572 (1985).
		\href{https://doi.org/10.1016/0550-3213(85)90002-1}{doi:10.1016/0550-3213(85)90002-1}.
		
		\bibitem[8]{Ruhe2023a}
		D.~Ruhe \textit{et al.},
		``Clifford group equivariant neural networks,''
		in \textit{Advances in Neural Information Processing Systems} (NeurIPS 2023).
		arXiv:2305.11141.
		
		\bibitem[9]{Brehmer2023}
		J.~Brehmer \textit{et al.},
		``Geometric algebra transformer,''
		in \textit{Advances in Neural Information Processing Systems} (NeurIPS 2023).
		arXiv:2305.18415.
		
		\bibitem[10]{Brehmer2024}
		L.~Brehmer \textit{et al.},
		``L-GATr: Lorentz-equivariant geometric algebra transformer,''
		(2024).
		arXiv:2405.14806.
		
		\bibitem[11]{Ruhe2024}
		D.~Ruhe \textit{et al.},
		``Clifford-steerable convolutional neural networks,''
		in \textit{International Conference on Machine Learning} (ICML 2024).
		arXiv:2402.14730.
		
		% NOTE: References [12]–[19] are reserved and not assigned in this paper.
		% This gap preserves continuity with the Furey et al.\ block numbered
		% [20]–[26] in Paper~1 (\cite{FirestoneP1}).
		
		%% [20]–[28]: Furey programme, division algebras, and software
		%% Numbers enforced via \bibitem[N]{key} to match Paper 1 reference list.
		\bibitem[20]{Furey2012}
		C.~Furey,
		``Unified theory of ideals,''
		\textit{Phys.\ Rev.}\ D~\textbf{86}, 025024 (2012).
		\href{https://doi.org/10.1103/PhysRevD.86.025024}{doi:10.1103/PhysRevD.86.025024}.
		
		\bibitem[21]{Furey2014}
		C.~Furey,
		``Generations: Three prints, in colour,''
		\textit{JHEP}~\textbf{2014}, 046 (2014).
		\href{https://doi.org/10.1007/JHEP10(2014)046}{doi:10.1007/JHEP10(2014)046}.
		
		\bibitem[22]{Furey2015}
		C.~Furey,
		``Standard model physics from an algebra?''
		Ph.D. thesis, University of Waterloo (2015).
		arXiv:1611.09182.
		
		\bibitem[23]{Furey2018a}
		C.~Furey,
		``Three generations, two unbroken gauge symmetries, and one
		eight-dimensional algebra,''
		\textit{Phys.\ Lett.}\ B~\textbf{785}, 84 (2018).
		\href{https://doi.org/10.1016/j.physletb.2018.08.032}{doi:10.1016/j.physletb.2018.08.032}.
		
		\bibitem[24]{Furey2018b}
		C.~Furey,
		``$\su(3)_C\times\su(2)_L\times\mathrm{U}(1)_Y$
		$(\times\,\mathrm{U}(1)_X)$ as a symmetry of division algebraic ladder
		operators,''
		\textit{Eur.\ Phys.\ J.}\ C~\textbf{78}, 375 (2018).
		\href{https://doi.org/10.1140/epjc/s10052-018-5844-7}{doi:10.1140/epjc/s10052-018-5844-7}.
		
		\bibitem[25]{Dixon1994}
		G.~M. Dixon,
		\textit{Division Algebras: Octonions, Quaternions, Complex Numbers and the
			Algebraic Design of Physics}
		(Kluwer Academic, Dordrecht, 1994).
		\href{https://doi.org/10.1007/978-1-4757-2315-1}{doi:10.1007/978-1-4757-2315-1}.
		
		\bibitem[26]{Stoica2018}
		O.~C. Stoica,
		``Leptons, quarks, and gauge from $C\ell_6$,''
		\textit{Adv.\ Appl.\ Clifford Algebras}~\textbf{28}, 52 (2018).
		\href{https://doi.org/10.1007/s00006-018-0869-4}{doi:10.1007/s00006-018-0869-4}.
		
		\bibitem[27]{Harris2020}
		C.~R. Harris \textit{et al.},
		``Array programming with {NumPy},''
		\textit{Nature}~\textbf{585}, 357 (2020).
		\href{https://doi.org/10.1038/s41586-020-2649-2}{doi:10.1038/s41586-020-2649-2}.
		
		\bibitem[28]{Virtanen2020}
		P.~Virtanen \textit{et al.},
		``{SciPy}~1.0: Fundamental algorithms for scientific computing in {Python},''
		\textit{Nature Methods}~\textbf{17}, 261 (2020).
		\href{https://doi.org/10.1038/s41592-019-0686-2}{doi:10.1038/s41592-019-0686-2}.
		
	\end{thebibliography}
	
\end{document}